\documentclass[11pt,a4paper]{article}
\usepackage[margin=2.9cm]{geometry}
\usepackage{amsmath,amssymb,amsthm,mathrsfs}
\usepackage[T1]{fontenc}
\usepackage{microtype}
\usepackage{booktabs}
\usepackage[colorlinks=true,linkcolor=blue,citecolor=blue]{hyperref}

% ------------------------------------------------------------------
% Part I unified manuscript.  Phase status: U4 complete (adversarial pass applied).
% Build: pdflatex x3.  Gates and step maps: see part_i/README.md and
% strategy.md, section "Part I unification - plan (recorded 2026-08-10)".
% Transcription discipline: displays copied from their source of truth
% carry a "% src:" comment.  Do not retype constants from memory.
% ------------------------------------------------------------------

\theoremstyle{plain}
\newtheorem{theorem}{Theorem}[section]
\newtheorem{proposition}[theorem]{Proposition}
\newtheorem{lemma}[theorem]{Lemma}
\newtheorem{corollary}[theorem]{Corollary}
\newtheorem{hypothesis}[theorem]{Hypothesis}
\theoremstyle{definition}
\newtheorem{definition}[theorem]{Definition}
\newtheorem{remark}[theorem]{Remark}
\newtheorem{srcfact}[theorem]{Source Fact}
\newtheorem*{srcfactstar}{Source Fact}
\newtheorem{assumption}{Assumption}
\renewcommand{\theassumption}{A\arabic{assumption}}

\newcommand{\R}{\mathbb{R}}
\newcommand{\Z}{\mathbb{Z}}
\newcommand{\C}{\mathbb{C}}
\newcommand{\E}{\mathbb{E}}
\newcommand{\A}{\mathfrak{A}}
\newcommand{\G}{\mathcal{G}}
\newcommand{\CE}{\mathcal{C}_{\mathrm E}}
\newcommand{\CH}{\mathcal{C}_{\mathrm H}}
\newcommand{\HH}{\mathcal{H}}
\newcommand{\Om}{\Omega}
\newcommand{\ip}[2]{\langle #1,#2\rangle}
\DeclareMathOperator{\supp}{supp}
\DeclareMathOperator{\dist}{dist}
\DeclareMathOperator{\spec}{spec}
\DeclareMathOperator{\Span}{span}
\DeclareMathOperator{\Ran}{Ran}

\title{Removal of the spatial cutoff for weakly coupled $P(\varphi)_2$:\\
norm convergence of the ground-state net on the quasi-local algebra\\[4pt]
\large Part I --- unified manuscript}
\author{Alexander Stottmeister\\[3pt]
  \normalsize Institut f\"ur Theoretische Physik, Leibniz Universit\"at Hannover,\\
  \normalsize Appelstra\ss e 2, 30167 Hannover, Germany}
% Fixed date on purpose: it tracks the last mathematical change (the scope
% amendment of A8(vi), 5 September 2026), not the build date, so the title page
% does not drift on every rebuild.  Update it only when the mathematics changes.
% The PDF bytes still carry pdftex's build timestamp; for byte-identical output
% build with SOURCE_DATE_EPOCH=1788566400 FORCE_SOURCE_DATE=1 (verified).
\date{5 September 2026}

\begin{document}
\maketitle

\begin{quote}\small
\textbf{Repair status (16 August 2026).} This is the canonical Part I manuscript.
The proof below is an implication from the imports in \S\ref{subsec:imports}; it is
not a proof of those imports.  In particular, the exact slab-uniform moment estimate
$U_{\rm GJS}$ is now stated explicitly in Assumption~\ref{A:gjs}(vi).  The printed
GJS text gives strong evidence for that reading, but Appendix~\ref{app:huniform} is a
source-mechanism audit and a proof of analytic prerequisites, not a reconstruction of
the GJS cluster expansion.  Consequently Theorem~\ref{thm:MAIN} is conditional on
$U_{\rm GJS}$ and the other enumerated imports.  The localization omission found in
the CE2-to-gap passage and the ultraviolet-mollifier specification are repaired below.

\smallskip
\textbf{Scope amendment (5 September 2026).}  The status of $U_{\rm GJS}$ is re-based.
It is no longer described as a quantifier strengthened beyond the source: it is the
cutoff-uniform quantifier that \cite{GJS} itself states (p.~629) and consumes in its
printed proofs of Theorems 1.1.7 and 1.1.1 (pp.~594--597).  Assumption~\ref{A:gjs}(vii$'$)
quotes the passages verbatim and the Scope paragraph of Assumption~\ref{A:gjs} records
the reading.  $U_{\rm GJS}$ remains a named hypothesis of Theorem~\ref{thm:MAIN}: it is
imported at the source's own scope, not re-derived here.  No proof in the manuscript
changed.
\end{quote}

\begin{abstract}
% src: m6/main.tex abstract (v2, audited repair), transcribed; forward pointers adapted.
Let $H(g)$ be the spatially cutoff $P(\varphi)_2$ Hamiltonian, $\Om_g$ its ground state, and
$\omega_g$ the induced state on the time-zero quasi-local algebra $\A$. Glimm and Jaffe (1970)
constructed an infinite-volume vacuum as a $w^*$ \emph{limit point} of space-averaged $\omega_g$;
the cited construction does not prove convergence of the net itself. Guerra--Rosen--Simon
(1975) formulated analogous Euclidean local-convergence questions as their Problems~1--2.
Assuming the exact slab-uniform interface $U_{\rm GJS}$ of
Assumption~\ref{A:gjs}(vi), we prove: in the Glimm--Jaffe--Spencer weak-coupling region
$0\le\lambda/m_0^2<\varepsilon_{\mathscr P}$, the \emph{full net} $(\omega_g)$ --- no subnets, no
averaging --- converges in norm on every local algebra,
\[
  \lim_{g\in\G}\bigl\|(\omega_g-\omega_\infty)|_{\A(O)}\bigr\|=0 ,
\]
with an explicit super-polynomial rate in $\dist(O,\{g<1\})$; the limit is locally normal,
translation invariant, $\Z_2$-invariant for even $\mathscr P$, and is a ground state of the
Glimm--Jaffe dynamics $(\A,\alpha)$ in the spectral sense. The proof is Hamiltonian: a
local-perturbations-perturb-locally estimate built on an exact spectral filter and the exact
light cone makes the net Cauchy; its two hypotheses --- a cutoff-uniform gap
$\spec H(g)\subseteq\{0\}\cup[m_1,\infty)$ and a uniform local interaction moment --- are derived
from the exponential-clustering theorem of Glimm--Jaffe--Spencer (Ann.\ of Math.\ 100
(1974), Thm.~1.1.7) and the working-family version of their moment estimate
(Thm.~1.1.8).  The latter is precisely $U_{\rm GJS}$: it is the cutoff-uniform
quantifier that the source itself states (p.~629) and consumes in its printed proofs of
Theorems 1.1.7 and 1.1.1 (pp.~594--597); it is imported here at that scope, not
re-proved.
The proof also imports, and explicitly identifies in \S\ref{sec:setting}, the standard
cutoff-construction results on self-adjointness and the simple positive ground state,
Feynman--Kac--Nelson transfer, finite-volume stability and hypercontractivity, domain smoothing
and a common core, finite propagation and patching, and covariance. No Euclidean uniqueness
theorem, global Markov property for the infinite-volume measure, or reverse
Osterwalder--Schrader reconstruction is used. We do not claim uniqueness of $\omega_\infty$
among all locally normal ground states, a sharp exponential rate, or anything outside the
weak-coupling region.
\end{abstract}

\tableofcontents

\section{Introduction}\label{sec:intro}

\subsection{The problem}

% src: m6/main.tex \S1.1 (v2), transcribed.
The $\lambda\varphi^4_2$ (more generally $P(\varphi)_2$) model was constructed by Glimm and
Jaffe in the series I--IV \cite{GJ1,GJ2,GJ3,GJ4}. Paper III builds the physical vacuum as
follows: the finite-volume ground states $\Om_g$ of the cutoff Hamiltonians $H(g)$ induce states
$\omega_g$ on the quasi-local algebra $\A$; after a space-translation averaging, a
$w^*$-compactness argument produces a \emph{limit point} $\omega$, which is shown to be locally
Fock \cite[Thms.~2.1--2.3]{GJ3}. Whether the net $(\omega_g)$ \emph{converges} --- so that the
vacuum is canonical rather than selected by a limit procedure --- was not proved there; the
Expositions volume states plainly that the $g\to1$ limit of the vacuum expectation values ``has
not been proved'' \cite[\S4]{GJExp}. Guerra, Rosen and Simon list analogous \emph{Euclidean}
local $L^p$ and finite-volume Gibbs-state convergence statements as Problems~1--2 of
\cite[pp.~125--126]{GRS75}; their note added on p.~128 records Newman's solution of the
small-coupling $p=1$ versions. (Some later surveys state the Hamiltonian cutoff limit as a
convergence; the cited primary construction itself is unambiguously a limit-point
construction.)

On the Euclidean side, convergence of \emph{Schwinger functions} at weak coupling follows from
the cluster expansion of Glimm--Jaffe--Spencer \cite{GJS}. Albeverio--H\o egh-Krohn--Zegarlinski
prove Euclidean Gibbs uniqueness under their stated hypotheses \cite{AHZ};
Bauerschmidt--Dagallier--Weber prove uniqueness of the invariant measure for the $\Phi^4_2$
stochastic-quantization equation above the critical temperature \cite{BDW}. Neither statement
is, by itself, a theorem about norm convergence of this Hamiltonian ground-state cutoff net.
The cited primary construction does not prove that local-norm convergence, and no such theorem
is identified in the references reviewed here. The Hamiltonian-renormalisation programme of
Rodriguez Zarate--Thiemann cited here is explicitly restricted to finite volume \cite{RZT}.

\subsection{The result}

\begin{quote}
\emph{Assuming $U_{\rm GJS}$, in the GJS weak-coupling region the full net
$(\omega_g)_{g\in\G}$ converges in norm on
every local algebra; the limit is locally normal, translation and (for even $\mathscr P$)
$\Z_2$ invariant, and is a ground state of the Glimm--Jaffe dynamics. An explicit rate holds:
$\|(\omega_g-\omega_\infty)|_{\A(O)}\|\le\Psi_{m_1}(d)$ for $g\equiv1$ on $O_d$, with
$\Psi_{m_1}$ decaying faster than every inverse power.}
\end{quote}

See Theorem~\ref{thm:MAIN}. The logical structure is deliberately Hamiltonian and
compactness-free:

\begin{center}
\begin{tabular}{ccc}
\toprule
step & content & where proved\\
\midrule
(i) & printed GJS CE2 and the source-scope CE1 interface $U_{\rm GJS}$
    & Assumption~\ref{A:gjs} (\S\ref{sec:setting})\\
(ii) & uniform local moment (M$_2$) & \S\ref{sec:M2}\\
(iii) & uniform gap (G): $\spec H(g)\subseteq\{0\}\cup[m_1,\infty)$ & \S\ref{sec:G}\\
(iv) & LPPL/Cauchy: (G)+(M$_2$) $\Rightarrow$ full-net convergence + rate & \S\S\ref{sec:path}--\ref{sec:limit}\\
(v) & the limit is an $\alpha$-ground state; invariances & \S\ref{sec:ground}\\
\bottomrule
\end{tabular}
\end{center}

Steps (ii)--(iii) convert Euclidean input to Hamiltonian statements through slab
Feynman--Kac--Nelson transfers in which every limit is taken at fixed $g$; no uniform gap is
presupposed anywhere in their proofs (the gap of step (iii) is an output). Step (iv) replaces
the compactness of \cite{GJ3} by a Cauchy argument: the same gap that yields clustering also
yields a local-perturbations-perturb-locally estimate, which bounds
$\|(\omega_g-\omega_{g'})|_{\A(O)}\|$ whenever both cutoffs equal $1$ on $O_d$. Step (v) is an
exact spectral identity: the ground-state property of the limit costs an identity, not an
estimate.

The dependency graph, established by the independent audit of the predecessor manuscript
(Appendix~\ref{app:provenance}), is:
% src: audit_2026/critical_proof_audit.tex \S3, transcribed.
\begingroup\small
\[
\begin{array}{c}
\text{cutoff self-adjointness, simple positive ground state, FKN, domains}\\
\quad+\quad\text{GJS CE2, }U_{\rm GJS}\text{ (slab CE1), scaling}
\end{array}
\Longrightarrow
\begin{array}{c}
(\mathrm M_2)\ \text{uniform local moment}\\
(\mathrm G)\ \text{uniform cutoff gap}
\end{array}
\]
\[
\begin{array}{c}
(\mathrm M_2)+(\mathrm G)+\text{finite propagation}+\text{patching}\\
+\text{affiliation}+\text{path regularity}
\end{array}
\Longrightarrow
\text{LPPL estimate}
\Longrightarrow
\begin{array}{c}
\text{full-net local-norm convergence}\\
\text{local normality}
\end{array}
\]
\[
\text{local-norm convergence}+\text{patching}+\text{covariance}
\Longrightarrow
\begin{array}{c}
\text{spectral ground state}\\
\text{translation and parity invariance.}
\end{array}
\]
\endgroup

\subsection{What is not claimed}

% src: m6/main.tex \S1.3 (v2), transcribed; refs adapted.
(a) \emph{Uniqueness}: Theorem~\ref{thm:MAIN} identifies the limit of \emph{this} cutoff family;
it does not show that every locally normal ground state of $(\A,\alpha)$ equals $\omega_\infty$.
(b) \emph{Sharp rate}: in the exactly solvable quadratic model the true rate is
$\mathrm{poly}(d)\,e^{-2m_1d}$ (companion note \texttt{m2/}); $\Psi_{m_1}$ is super-polynomial but not
exponential, and this is intrinsic to the filter method used. (c) \emph{Beyond weak coupling}:
under the explicit phase-classification hypotheses stated in Hypothesis~\ref{hyp:broken}, a
cutoff-uniform gap is incompatible with that two-phase structure
(Theorem~\ref{thm:dichotomy}). Thus an extension to such a broken phase would have to abandon
the uniform-gap input; no claim is made outside $0\le\lambda/m_0^2<\varepsilon_{\mathscr P}$.
(d) \emph{Wavelet uniformity}: uniformity in a lattice/wavelet resolution is a separate
problem (Part II of this programme) and is not addressed here.

\subsection{The proof strategy in operator-algebraic terms}\label{subsec:strategy}

% Expository subsection, added after gate G6 at the author's request. It introduces no
% new mathematical input: every assertion is a cross-reference into the body, or is
% explicitly contextual.

\emph{The topology, and why it is the right one.} Each $\omega_g$ is the vector state of
the simple positive ground vector $\Om_g$ of $H(g)$, restricted to the quasi-local
algebra $\A$; the cutoffs form a directed set under inclusion of plateaus. The classical
construction \cite{GJ3} produced a vacuum as a $w^*$ \emph{limit point} of space-averaged
$\omega_g$: compactness selects a state but proves neither convergence nor canonicity.
Theorem~\ref{thm:MAIN} asserts convergence of the full, unaveraged net in the norm of
every local dual $\A(O)^*$ --- the strongest locally uniform statement available --- and
Corollary~\ref{cor:canonical} then makes the historical construction subsequence- and
averaging-independent. Norm convergence on local duals also transports normality: the
limit lies in the local Fock folia (Proposition~\ref{prop:normal}).

\emph{Why standard tools fail.} The difference $V(g')-V(g)$ has infinite norm, so the
quasi-adiabatic/automorphic-equivalence machinery for gapped lattice systems
\cite{Hastings,BMNS} does not start; ergodicity-type arguments give extremality of limit
points, not uniqueness; and Euclidean convergence of Schwinger functions lives in the
wrong topology --- there is no soft passage from moments to norms on local von Neumann
algebras. The architecture here is instead a \emph{Cauchy} argument built from a
local-perturbations-perturb-locally estimate: if $g\equiv g'\equiv1$ on $O_d$, then
$\|(\omega_g-\omega_{g'})|_{\A(O)}\|\le\Psi_{m_1}(d)$ with $\Psi_{m_1}$ super-polynomial
(Proposition~\ref{prop:onepath}, Lemma~\ref{lem:rate}). Plateaus are cofinal and local
duals are complete, so the net converges (Theorem~\ref{thm:cauchy}) --- no compactness
anywhere.

\emph{Pillar 1: two cutoff-uniform Hamiltonian estimates from Euclidean input.} The
Cauchy estimate consumes exactly two nonstandard inputs, both uniform over the whole
Hamiltonian cutoff class $\CH$: the gap \eqref{eq:G} of Theorem~\ref{thm:G}, and the
local moment \eqref{eq:M2} of Theorem~\ref{thm:M2}. Both are transferred from the two
cluster-expansion estimates of \cite{GJS} for the interacting Euclidean measures
$dq_h$ --- the moment bound of Source Fact~\ref{sf:CE1}, used through the
working-family interface $U_{\rm GJS}$ (imported at the scope at which the source
states and uses it, Assumption~\ref{A:gjs}(vii$'$)), and the exponential clustering of
Source Fact~\ref{sf:CE2}.  Appendix~\ref{app:huniform} records the source mechanism and proves
the normalization and stability prerequisites; it does not replace the source's
polymer expansion. The transfer device is the
Feynman--Kac--Nelson slab formula: $\psi_T=e^{-TH(g)}\Om_0$ normalizes to the ground
state, and sharp-time Euclidean moments become matrix elements of $e^{-tH(g)}$. One
structural point deserves emphasis: $\hat\psi_T\to\Om_g$ uses \emph{simplicity only}
(Lemma~\ref{lem:slabvac}) --- no gap is presupposed, and the gap of
Theorem~\ref{thm:G} is an output, so there is no circularity. Converting semigroup decay
on centered polynomial vectors into a spectral gap requires that these vectors be total;
this is the density theorem (Theorem~\ref{thm:dense}), proved from the factorial moments
of Source Fact~\ref{sf:CE1} via exponential integrability, strip analyticity of
characteristic functionals, Fourier uniqueness, and martingale convergence.

\emph{Pillar 2: the derivative estimate.} Along the affine interpolation
$g_s=(1-s)g+sg'$, first-order perturbation theory for the simple ground state gives
$\dot\Om_s=-H_s^{-1}Q_sV(v)\Om_s$ (Proposition~\ref{prop:projreg}). Three moves control
this, each with an operator-algebraic flavor. \emph{(i) An exact inverse-energy filter}
(Lemma~\ref{lem:filter}) writes $H^{-1}(1-P_0)=\int W_\gamma(t)\,e^{itH}(1-P_0)\,dt$,
converting the reduced resolvent into dynamics. The kernel $W_\gamma$ decays faster than
every inverse power but \emph{cannot} decay exponentially: a kernel with exponential
tails has a transform analytic in a strip, which cannot vanish on a real neighbourhood of
$0$ --- as $\widehat W_\gamma$ must --- without vanishing identically. This is why the
final rate is super-polynomial rather than exponential; it is intrinsic to the method
(cf.\ \S\ref{sec:discussion}). \emph{(ii) The exact light cone}: $e^{itH_s}$ turns
$A\in\A(O)$ into $\beta^{g_s}_t(A)\in\A(O_{|t|})$ with speed $\le1$ exactly
(Proposition~\ref{prop:propagation}); no Lieb--Robinson theory is needed. \emph{(iii)
Clustering with one unbounded local argument} (Lemma~\ref{lem:cluster}): the correlation
between a far perturbation cell $V_j$ and the localized $\beta_t(A)$ is bounded by
$2\|V_j\Om_s\|\|A\|e^{-\gamma R/2}$. The mechanism is a quantitative substitute for the
strip-analyticity arguments familiar from positive-energy representations: the two
boundary functions have spectral supports in $[\gamma,\infty)$ and $(-\infty,-\gamma]$,
locality makes them agree on the interval $|t|<R$, and an optimized Gaussian separator
(Lemma~\ref{lem:separation}) bounds the value at $t=0$. Note what the moment
\eqref{eq:M2} is doing: the perturbation is unbounded, but only $V_j$ \emph{evaluated in
the ground state} enters --- the uniform moment is the operator-algebraic softening of
the unboundedness. A unit-scale partition of the perturbation, the two-regime split of
each time integral, and the count of at most eight cells per distance shell assemble the
estimate (Lemma~\ref{lem:statederiv} through Proposition~\ref{prop:onepath}).

\emph{Assembly and identification.} Cauchy-ness in each $\A(O)^*$ gives compatible limit
functionals and the state $\omega_\infty$; local normality follows by a monotone-net
argument (Proposition~\ref{prop:normal}). The ground-state property costs an identity,
not an estimate: for each cutoff and $|t|<d(g,O)$, patching \eqref{eq:P} makes
$\omega_g(A\alpha_t(B))$ literally the Fourier--Stieltjes transform of a measure
supported in $[0,\infty)$, and the positive-energy condition passes to the limit by
locally uniform convergence (Proposition~\ref{prop:spectral}). Translation invariance
needs no averaging: $g\mapsto\tau_ag$ is an order isomorphism of the cutoff net with
cofinal range, so covariance \eqref{eq:T} plus net convergence suffices; parity likewise.

\emph{Boundaries.} All of this sits in the weak-coupling region through the explicit
import ledger of \S\ref{subsec:imports} and the coordinate convention of
Remark~\ref{rem:wlog}. The theorem does not identify $\omega_\infty$ among \emph{all}
locally normal ground states of $(\A,\alpha)$, does not reach a sharp exponential rate
(the filter obstruction above; the solvable quadratic model has
$\mathrm{poly}(d)\,e^{-2m_1d}$), and claims nothing beyond weak coupling; there,
Theorem~\ref{thm:dichotomy} records the expected mechanism --- for even $\mathscr P$, a
cutoff-uniform gap is incompatible with the fully classified two-phase structure of
Hypothesis~\ref{hyp:broken}. In one sentence: \emph{the same uniform gap that gives
exponential clustering also gives a local-perturbation stability estimate through an
exact inverse-energy filter and the exact light cone; the cluster expansion of
\cite{GJS} supplies the gap and one local moment uniformly in the cutoff; and
Cauchy-ness in the local duals replaces every compactness step of the classical
construction.}

\subsection*{Provenance and relation to the frozen sources}

This document unifies, at full explicitness, the following frozen components of the research
record (file pointers in Appendix~\ref{app:provenance}): the audited manuscript
\texttt{m6/main.tex} (v2), the two structured proof reconstructions
\texttt{audit\_2026/}\allowbreak\texttt{lamport\_euclidean\_to\_gap.tex} and
\texttt{audit\_2026/}\allowbreak\texttt{lamport\_gap\_to\_main.tex}, the source-scope and constant gates
\texttt{m4/r1-check.tex} and \texttt{m4/r2-check.tex}, the spectral-passage note \texttt{m0/},
and the dichotomy note \texttt{m3/}. Every repair identified by the independent audit of
9~August 2026 (register in Appendix~\ref{app:provenance}) is incorporated. External inputs are
the printed theorems of \cite{GJS} and the explicitly cited construction results from
\cite{GJ1,GJ2,GJ71,Rosen70,SimonBook,GRS75,GJExp}; they are enumerated once, as
Assumptions~\ref{A:cutoff}--\ref{A:gjs}, and no other external mathematical input is
used (Corollary~\ref{cor:canonical} additionally quotes the averaged-state formula
\cite[(2.3)--(2.5)]{GJ3} as the definition of the averaged states it identifies).

\section{Setting, notation, and standing imports}\label{sec:setting}

\subsection{The model}\label{subsec:model}

% src: m6/main.tex \S2.1 (v2), transcribed.
$\HH=L^2(\mathcal S'(\R),d\phi_0)$ is $Q$-space over the time-zero free field of mass $m_0>0$:
$d\phi_0$ is Gaussian with covariance $C=(2\mu)^{-1}$, $\mu=(-\partial_x^2+m_0^2)^{1/2}$. Weyl
operators $W(f,h)=\exp i(\varphi_0(f)+\pi_0(h))$ generate the local net
\[
  \A(O)=\{W(f,h):f,h\in C_c^\infty(O)\ \text{real}\}'',\qquad
  \A=\overline{\textstyle\bigcup_O\A(O)}^{\;\|\cdot\|} ,
\]
$O$ ranging over bounded open intervals. Let
$\mathscr P(\xi)=\sum_{n=0}^{2p'}a_n\xi^n$, $a_{2p'}>0$, $p'\ge1$, and for real
$h\in L^1(\R)\cap L^2(\R)$ of compact support
\[
  V(h):=\lambda\int h(x)\,{:}\mathscr P(\varphi_0(x)){:}_C\,dx ,
\]
a self-adjoint multiplication operator in $Q$-space, affiliated with $\A(U)$ for every bounded
open $U\supseteq\supp h$, and finitely additive on common domains
($V(\sum_j h_j)\xi=\sum_j V(h_j)\xi$ whenever all terms are defined). For a cutoff $g$ as
below,
\[
  K(g):=H_0+V(g),\quad E(g):=\inf\spec K(g),\quad H(g):=K(g)-E(g),
\]
with normalized ground state $\Om_g$ (Assumption~\ref{A:cutoff}),
$\omega_g:=\ip{\Om_g}{\cdot\,\Om_g}$, and $\beta^g_t=\mathrm{Ad}\,e^{itH(g)}$.

\subsection{The three cutoff classes}\label{subsec:classes}

\begin{definition}[cutoff classes]\label{def:classes}
\begin{align*}
 \CE&:=\{h:\R^2\to[0,1]\ \text{measurable},\ \supp h\ \text{compact}\}
 &&\text{(Euclidean class)},\\
 \CH&:=\{g:\R\to[0,1]\ \text{measurable},\ \supp g\ \text{compact}\}
 &&\text{(Hamiltonian class)},\\
 \G&:=\{g\in C_c^\infty(\R):0\le g\le1,\ g\equiv1\ \text{on a nonempty open interval}\}
 &&\text{(plateau net class)}.
\end{align*}
The first two are the printed classes of \cite[p.~589]{GJS} and
\cite[(1.6), p.~588]{GJS}, respectively.
$\G$ is directed by $g\preccurlyeq g'\iff\{g=1\}^\circ\subseteq\{g'=1\}^\circ$. For a bounded
interval $O$ and $r\ge0$ set
\[
 O_0:=O,\qquad O_r:=\{x\in\R:\dist(x,O)<r\}\quad(r>0),\qquad
 d(g,O):=\dist\bigl(O,\R\setminus\{g=1\}^\circ\bigr).
\]
\end{definition}

The roles are fixed once and for all: the Euclidean estimates of Assumption~\ref{A:gjs}
quantify over $\CE$; the uniform moment and gap of \S\S\ref{sec:M2}, \ref{sec:G} quantify over
$\CH$; the convergence theorem quantifies over the subnet class $\G\subset\CH$. The elementary
compatibility statements (memberships and affine-path stability) are
Lemma~\ref{lem:classes}.

\subsection{Notation freeze}\label{subsec:freeze}

All later sections conform to the following table; in particular every symbol imported from a
frozen source is converted to this notation at transcription time (see
\texttt{part\_i/README.md} for the conversion log).

\begin{center}
\small
\begin{tabular}{l p{0.66\textwidth}}
\toprule
symbol & meaning / convention\\
\midrule
$\mathscr P$, $2p'$ & interaction polynomial; $\deg\mathscr P=2p'$, $a_{2p'}>0$, $p'\ge1$\\
$\ip{\cdot}{\cdot}$ & inner product, conjugate-linear in the \emph{first} argument\\
$C$, $\mu$ & time-zero covariance $C=(2\mu)^{-1}$, $\mu=(-\partial_x^2+m_0^2)^{1/2}$\\
$\CE,\ \CH,\ \G$ & cutoff classes of Definition~\ref{def:classes}, in the roles fixed there\\
$K(g),E(g),H(g),\Om_g,\omega_g,\beta^g$ & cutoff Hamiltonian data (\S\ref{subsec:model})\\
$\alpha$, $\tau_a$, $U_a$ & patched dynamics (Assumption~\ref{A:patch}); space translations and their $Q$-space unitaries\\
$\Theta$, $\theta$ & parity unitary $\Theta=\Gamma(-1)$; parity automorphism $\theta=\mathrm{Ad}\,\Theta$\\
$K_1,K_2,C_7,C_8,m_1$ & GJS constants (Assumption~\ref{A:gjs}); normalization $K_2\ge1$\\
$\|\cdot\|_{\rm GJS}$ & printed observable norm (1.1.7)--(1.1.9) of \cite{GJS}\\
$N_I$ & number of unit cells $[j,j+1)$ meeting the interval $I$ in positive measure\\
$M$ & uniform local moment constant of (M$_2$) (\S\ref{sec:M2})\\
$\gamma$ vs.\ $m_1$ & $\gamma$ = abstract gap parameter in \S\S\ref{sec:path}--\ref{sec:lppl}; instantiated $\gamma=m_1$ in \S\ref{sec:limit}\\
$\mathfrak m$ & total-variation bound in the spectral-separation lemma (\S\ref{sec:filter})\\
$h_{T,g},\ \psi_s,\ \hat\psi_s,\ c_g$ & slab objects (Assumption~\ref{A:fkn}; \S\ref{sec:slab})\\
$\Psi_\gamma$ & rate function of the one-path estimate (\S\ref{sec:lppl})\\
$d_0$ & LPPL threshold $d_0=4+2/\gamma$\\
\bottomrule
\end{tabular}
\end{center}

Two deliberate renamings relative to the frozen sources: the separation-lemma variation bound
is $\mathfrak m$ (it was $M_0$ in \texttt{lamport\_gap\_to\_main.tex}, colliding with the
moment constant $M$), and $\Theta=\Gamma(-1)$ is reserved for the parity unitary throughout.

\subsection{Standing imports}\label{subsec:imports}

Each import below states (a) the exact statement used, (b) the pinpoint source, (c) a scope
remark verifying that the used statement is no stronger than the printed one, (d) the sections
that consume it, and (e) the page images documenting it in the audit dossier
(\texttt{audit\_2026/}\allowbreak\texttt{evidence/}, index in
\texttt{audit\_2026/}\allowbreak\texttt{citation\_screenshot\_dossier.tex}).

\begin{assumption}[cutoff construction]\label{A:cutoff}
For every $g\in\CH$: $K(g)=H_0+V(g)$ is essentially self-adjoint and bounded below; its closure
(again denoted $K(g)$) has $E(g)=\inf\spec K(g)$ as a \emph{simple} eigenvalue, and the
normalized eigenvector $\Om_g$ may be chosen strictly positive $d\phi_0$-a.e.\ in $Q$-space.
For every real $h\in L^1\cap L^2$ of compact support, $V(h)$ is self-adjoint on its maximal
multiplication domain.

\emph{Sources.} \cite[p.~588]{GJS}; existence and uniqueness of the cutoff vacuum:
\cite[Thms.~2.2.1, 2.3.1]{GJ2}; lower bounds and positivity: \cite{GJ71}; the semigroup
construction: \cite[Thm.~5.3]{Rosen70}.

\emph{Scope.} Strict positivity of $\Om_g$ is the standard positivity-improving semigroup
argument of the cited construction; simplicity is used below only through the two consequences
$c_g:=\ip{\Om_g}{\Om_0}>0$ and the absence of a second ground vector. \emph{Isolation of
$E(g)$ (a gap) is not imported}; the gap of \S\ref{sec:G} is an output.

\emph{Used in.} \S\ref{sec:slab} (slab limits), \S\ref{sec:M2} (weak-graph passage),
\S\ref{sec:G} (density, spectral projections), \S\ref{sec:path} (simple Perron root),
\S\ref{sec:ground} and \S\ref{sec:dichotomy} (covariance and parity via
Lemma~\ref{lem:covparity}).

{\raggedright\emph{Dossier.} \texttt{gj2-ground-existence-08}, \texttt{gjii-ground-12},
\texttt{gj71-lower-bound-02}, \texttt{gj71-theorem2-04}, \texttt{rosen-semigroup-22},
\texttt{rosen-selfadjoint-23}, \texttt{gjs-source-05/06} (E04--E06, E10--E13, E20).\par}
\end{assumption}

\begin{assumption}[Feynman--Kac--Nelson transfer]\label{A:fkn}
Let $d\mu_0$ be the free Euclidean Gaussian measure with covariance
$(-\Delta_{\R^2}+m_0^2)^{-1}$, $J_t$ the sharp-time embeddings, and, for $h\in\CE$,
$\mathcal U_h:=\lambda\int h(z){:}\mathscr P(\Phi(z)){:}\,dz$. For all bounded time-zero
functions $f,h'$, $T>0$, and $g\in\CH$,
\[
 \ip{f}{e^{-TK(g)}h'}
 =\int\overline{(J_Tf)}\,(J_0h')\,e^{-\mathcal U_{\mathbf 1_{[0,T]}\otimes g}}\,d\mu_0 ,
\]
with the interaction over exactly the slab $[0,T]$ carrying the two embeddings, together
with the analogous multi-time formulas on finite slabs (used in \S\ref{sec:slab} in the
two- and one-time forms, where the symmetric slab cutoffs
$h_{T,g}:=\mathbf 1_{[-T,T]}\otimes g\in\CE$ enter through the \emph{normalized}
measures $dq_{h_{T,g}}$).

\emph{Sources.} Free-field FKN: \cite[Thm.~III.6]{SimonBook}; interacting form:
\cite[Cor.~V.14, Thm.~V.15]{SimonBook}; equivalently \cite[Thms.~II.14, II.16]{GRS75}.

\emph{Scope.} \cite[Thm.~III.12]{SimonBook} is a projection-product estimate, % negative-scope mention (whitelisted): III.12
\emph{not} an FKN theorem, and is nowhere used (audit finding S1).
The identification of sharp-time Wick functionals with $Q$-space multiplication operators is
the conventions lock, Remark~\ref{rem:conv}.

\emph{Used in.} \S\ref{sec:slab} (slab identities), \S\ref{sec:path} (path formula).

{\raggedright\emph{Dossier.} \texttt{simon-free-fkn-111}, \texttt{simon-fkn-184/185},
\texttt{grs-fkn-39/40}, \texttt{simon-wrong-iii12-119} (E27--E43 range).\par}
\end{assumption}

\begin{assumption}[finite-volume stability]\label{A:stability}
For every $h\in\CE$: $\mathcal U_h\in L^p(d\mu_0)$ for every $p<\infty$, and
$e^{-\varrho\,\mathcal U_h}\in L^1(d\mu_0)$ for every finite $\varrho>0$. In particular
$0<Z(h):=\int e^{-\mathcal U_h}d\mu_0<\infty$, and the interacting measures
\[
 dq_h:=Z(h)^{-1}e^{-\mathcal U_h}\,d\mu_0 ,\qquad h\in\CE ,
\]
are well-defined probability measures.

\emph{Source.} \cite[Thm.~V.7]{SimonBook} and its chapter context.

\emph{Scope.} Only finite slabs and compactly supported $h$ occur below; no infinite-volume
stability is imported. A self-contained quantitative lower bound on $Z(h)$ and the
$h$-uniform $L^p$ bound on $e^{-\mathcal U_h}$ are re-proved in Appendix~\ref{app:huniform}.

\emph{Used in.} \S\ref{sec:reduction} (Lemma~\ref{lem:classes}), \S\ref{sec:slab},
\S\ref{sec:path} (H\"older/Vitali chains), Appendix~\ref{app:huniform}.

{\raggedright\emph{Dossier.} \texttt{simon-exp-174}, \texttt{simon-transfer-178/179/180} (E33--E43 range).\par}
\end{assumption}

\begin{assumption}[free hypercontractivity]\label{A:hyper}
The free time-zero Markov semigroup $\mathsf P_T$ satisfies
\[
 \|\mathsf P_TF\|_{q}\le\|F\|_{p}
 \qquad\text{whenever}\qquad \frac{q-1}{p-1}=e^{2m_0T},\quad 1<p\le q<\infty .
\]

More generally (the abstract form of the same source): for the second quantization
$\Gamma(c)$ of a contraction $c$ on the one-particle space of any Gaussian $L^2$-space,
$\|\Gamma(c)F\|_{q}\le\|F\|_{p}$ whenever $\|c\|^2\le(p-1)/(q-1)$; the display is the
case $c=e^{-T\mu}$, and the Mehler case $c=e^{-t}\mathbf 1$ on $L^2(d\mu_0)$ is the form
used in Appendix~\ref{app:huniform}.

\emph{Sources.} \cite[Thm.~I.17]{SimonBook} (abstract form and the display); time-region
form \cite[Thm.~III.17]{SimonBook}.

\emph{Scope.} \cite[Thms.~V.10, V.12]{SimonBook} are, respectively, a transfer-matrix % negative-scope mention (whitelisted): V.10, V.12
application and an essential-self-adjointness theorem; neither is the hypercontractivity input
and neither is cited for it (audit finding S2).

\emph{Used in.} \S\ref{sec:path} (the pairing bound), Appendix~\ref{app:huniform} (tail
estimate).

{\raggedright\emph{Dossier.} \texttt{simon-hyper-056}, \texttt{simon-timeslice-125/126}.\par}
\end{assumption}

\begin{assumption}[domain smoothing; common core]\label{A:domain}
\emph{(i)} For every $g\in\CH$ and $t>0$,
\[
 \Ran e^{-tK(g)}\subseteq D(H_0)\cap D(V')
\]
for every polynomial multiplication operator $V'$ of the time-zero field with coefficients in
$L^1\cap L^2$ of compact support (in particular for every $V(h)$ as in
\S\ref{subsec:model}).
\emph{(ii)} For every affine path $g_s=(1-s)g+sg'$, $g,g'\in\G$, there is a common core
$\mathscr C$ for all $K(g_s)$, $s\in[0,1]$, on which
$K(g_s)\xi=K(g_{s_0})\xi+(s-s_0)V(g'-g)\xi$.

\emph{Sources.} (i) is the model-specific smoothing theorem \cite[Ch.~7, Thm.~7.4]{GJExp}.
(ii) is part of the cutoff construction; essential self-adjointness pinpoints, in their
roles per the audit: \cite[Thm.~V.12]{SimonBook}, \cite[Thm.~5.3]{Rosen70}, \cite{GJ2}. % whitelisted: V.12 in its audited positive role (esa), not the S2 misuse

\emph{Scope.} \cite[Thm.~7.3]{GJExp} alone assumes a potential bounded below % negative-scope mention (whitelisted): Thm. 7.3
and is therefore \emph{not} sufficient here; the citation is to Theorem~7.4
(audit finding S3).

\emph{Used in.} Proposition~\ref{prop:smoothing}, \S\ref{sec:slab} (dominators),
\S\ref{sec:path} (eigenvalue differentiation).

{\raggedright\emph{Dossier.} \texttt{gjexp-domain-165/166/167} (E17--E19).\par}
\end{assumption}

\begin{assumption}[free propagation; patching; the dynamics $\alpha$]\label{A:patch}
\emph{(i)} The free dynamics $\mathrm{Ad}\,e^{itH_0}$ has propagation speed one
(Klein--Gordon finite propagation), and the Trotter product formula for $e^{itK(g)}$ converges
strongly.
\emph{(ii)} There is a one-parameter automorphism group $\alpha$ of $\A$, independent of $g$,
with
\begin{equation}
 \alpha_t(\A(O))\subseteq\A(O_{|t|}),\qquad
 \alpha_t(A)=\beta^g_t(A)\ \ \text{if } g\equiv1 \text{ on } O_{|t|}\ (A\in\A(O)).
 \tag{P}\label{eq:P}
\end{equation}

\emph{Sources.} \cite[\S III]{GJ1}; \cite[\S\S3.5--3.6]{GJ2}; \cite[\S4.1 and Ch.~4]{GJExp}.

\emph{Scope.} \eqref{eq:P} is the Glimm--Jaffe patching construction verbatim; the
$g$-uniform propagation bound for the \emph{cutoff} dynamics is not imported but derived
(Proposition~\ref{prop:propagation}).

\emph{Used in.} Proposition~\ref{prop:propagation}, \S\ref{sec:filter} (locality identity),
\S\ref{sec:lppl}, \S\ref{sec:ground}.

{\raggedright\emph{Dossier.} \texttt{gj1-locality-06}, \texttt{gjexp-054/055/056},
\texttt{gj2-cutoff-independence-38}, \texttt{gj2-local-algebra-40} (E03, E06,
E14--E16).\par}
\end{assumption}

\begin{assumption}[covariance data]\label{A:covar}
Let $U_a$ implement the space translation $\tau_a$ on $Q$-space and let
$\Theta:=\Gamma(-1)$, so that $(\Theta\psi)(q)=\psi(-q)$ in $Q$-space. Then, for every
$g\in\CH$,
\[
 U_aK(g)U_a^*=K(\tau_ag),\qquad \Theta H_0\Theta=H_0,\qquad
 \Theta V(g)\Theta=V(g)\ \ \text{for even }\mathscr P .
\]

\emph{Sources.} \cite[p.~398]{GJ2} (cutoff covariance); the transformation of Wick powers
under $\Gamma(-1)$ is the parity of the polynomial.

\emph{Scope.} Only the displayed operator identities are imported; the state-level
consequences are derived in Lemma~\ref{lem:covparity}.

\emph{Used in.} Lemma~\ref{lem:covparity}; thence \S\ref{sec:ground},
\S\ref{sec:dichotomy}.

{\raggedright\emph{Dossier.} \texttt{gj2-cutoff-independence-38} (E06).\par}
\end{assumption}

\begin{assumption}[the GJS package]\label{A:gjs}
The Euclidean input is \cite{GJS}, used exactly as printed, as follows.

\emph{(i) Measures.} For $h\in\CE$, $dq_h=Z(h)^{-1}e^{-\mathcal U_h}d\mu_0$ as in
Assumption~\ref{A:stability}; the Hamiltonian cutoff class of \cite[(1.6), p.~588]{GJS} is
$\CH$.

\emph{(ii) Observable class.} The class (1.1.5) of \cite{GJS} consists of sharp-time monomials
$Q(x^0)=\int d\vec x\,w(\vec x)\prod_i{:}\Phi(x_i)^{m_i}{:}$ with $L^2$ kernels; equal times
are allowed, and products of individually Wick-ordered powers are admissible without
reordering. GJS bound the degree of each \emph{individual} Wick factor: $n_i<\bar n$ for a fixed
finite $\bar n$, any choice $\bar n\ge\deg\mathscr P$ being sufficient; the printed text
adds ``Since there is no restriction on $r$, $Q$ still may have an arbitrarily large
degree'' (p.~593). Footnote~9 permits $\bar n=\infty$ at the price
$N(\Delta)!\to(N(\Delta)!)^3$ in (1.1.7), (1.1.8), (2.4.5). Every observable used below
has individual factor degrees $\le\max(2p',1)$, so the fixed finite choice
$\bar n=2p'+1$ applies and the printed first-power weights are used as displayed; the
\emph{total} degrees of the observables of \S\ref{sec:G} are unbounded, exactly as the
printed sentence licenses.

\emph{(iii) Printed norms.}
% src: m4/r2-check.tex, Source Fact 1.1 (page image of p.594), transcribed.
\begin{equation}
\begin{gathered}
 \textup{(1.1.7)}\ \ \|Q\|_{\rm GJS}=K_1\Bigl[\prod_\Delta K_2^{2N(\Delta)}N(\Delta)!\Bigr]\|w\|_2,\\
 \textup{(1.1.8)}\ \ \|Q(x_0)\|_{\rm GJS}=K_1\Bigl[\prod_\Delta K_2^{2N(\Delta)}N(\Delta)!\Bigr]\|w(x_0,\cdot)\|_2,
\end{gathered}
 \label{eq:printednorm}
\end{equation}
where $N(\Delta)$ is the total degree localized in the unit square $\Delta$, and
\textup{(1.1.9)}: a nonlocalized polynomial is normed by summing its localized pieces.
Replacing $K_2$ by $\max\{1,K_2\}$ only enlarges right-hand sides; \emph{we fix $K_2\ge1$}.

\emph{(iv) CE2.}
% src: m6/main.tex sf:CE2 (v2), transcribed.
\begin{srcfact}[CE2 $=$ {\cite[Thm.~1.1.7]{GJS}}, p.~594]\label{sf:CE2}
For localized monomials $Q_1,Q_2$ with disjoint localization squares separated by a strip of
width $d$ bounded by two parallel lines (any orientation),
\[
  \Bigl|\int Q_1Q_2\,dq_h-\int Q_1dq_h\int Q_2dq_h\Bigr|\le C_7e^{-m_1d}\|Q_1\|_{\rm GJS}\|Q_2\|_{\rm GJS} ,
\]
with $m_1>0$ and $C_7$ \emph{independent of $h$, $Q_1$, $Q_2$} (so stated in the source).
\end{srcfact}

\emph{(v) CE1.}
% src: m6/main.tex sf:CE1 (v2), transcribed.
\begin{srcfact}[CE1 $=$ {\cite[Thm.~1.1.8]{GJS}}, p.~595]\label{sf:CE1}
The printed estimate is
\[
 \bigl|\int Q(x^0)\,dq_h\bigr|\le C_8\|Q(x^0)\|_{\rm GJS},
\]
stated uniformly as the source cutoff tends to one.  This item records the literal
one-sentence theorem statement; the class-uniform quantifier actually consumed below,
which the source states on p.~629 and uses on pp.~595--597 (item~(vii$'$)), is isolated
in item~(vi) rather than being silently inserted here.
\end{srcfact}

\emph{(vi) Exact working-family interface $U_{\rm GJS}$ (imported at source scope).}
For
\[
 \CE^{\rm slab}:=
 \{h_{T,g}=\mathbf1_{[-T,T]}\otimes g:T>0,\ g\in\CH\}\subset\CE,
\]
one and the same constant $C_8$ satisfies
\begin{equation}
 \left|\int Q(x^0)\,dq_{h_{T,g}}\right|
 \le C_8\|Q(x^0)\|_{\rm GJS}
 \quad(T>0,\ g\in\CH) .
 \tag{$U_{\rm GJS}$}\label{eq:UGJS}
\end{equation}
The quantifiers include every affine cutoff
$g_s=(1-s)g+sg'$ used later, because $g_s\in\CH$.  Equation
\eqref{eq:UGJS}, and no uniformity assertion for a larger class, is the CE1
uniformity needed in \S\S\ref{sec:M2} and \ref{sec:G}.

\emph{(vii) Printed evidence for (vi).}
% src: m4/r1-check.tex sf:printed (page image of p.629), transcribed.
\begin{srcfact}[{\cite[p.~629, \S4]{GJS}}]\label{sf:printed}
``It is important to work with a translation invariant Hamiltonian, in order that time
translation preserve subspaces of bounded momentum. The estimates of Theorems 1.1.8, 1.1.11,
3.1 and 4.1 are uniform in the space cutoff $h$, and the path space integrals converge as
$h\to1$. Thus these theorems hold in the infinite volume limit, and we proceed to work in the
limit $h=1$.''
\end{srcfact}

\emph{(vii$'$) Printed use of the quantifier by the source.}  The following passages
show the source stating the cutoff-uniform quantifier of \eqref{eq:UGJS} and consuming
it for arbitrary cutoff functions of its class, not only in the limit $h\to1$.  The
source's $V(g)$ is the interaction $\mathcal U_g$ of item~(i); its
$\langle F\rangle_h$ is $\int F\,dq_h$.
% src: audit_2026/evidence/gjs-theorems-11.png (p.594), gjs-theorems-12.png (p.595),
% src: gjs-proof-14.png (p.597), gjs-uniform-46.png (p.629); transcribed from the page images.
\begin{srcfactstar}[{\cite[pp.~594--597 and 629]{GJS}}]
\emph{(a) Proof of Theorem 1.1.1, pp.~594--595.}  ``For test functions
$f_i\in C_0^\infty(\mathcal O)$ and for cutoff functions $h_1$ and $h_2$, we assert that
\[
 \Bigl|\int\prod_i\Phi(f_i)\,dq_{h_1}-\int\prod_i\Phi(f_i)\,dq_{h_2}\Bigr|
 \le O(1)e^{-m_1d/2}
 \tag{1.1.11}
\]
where $d=\operatorname{dist}(\mathcal O,\operatorname{suppt}\,h_1-h_2)$, and then
Theorem 1.1.1 follows.  To prove (1.1.11) let $g=h_2-h_1$, and define
$g_i=g\chi_{\Delta_i}$, where $\chi_{\Delta_i}$ is the characteristic function of
$\Delta_i$.  Also we use the notation $\langle F\rangle_h=\int F\,dq_h$, and set
$Q=\prod_i\Phi(f_i)$.  Then
\begin{align*}
 \langle Q\rangle_{h_2}-\langle Q\rangle_{h_1}
 &=\int_0^1d\alpha\,\frac{d}{d\alpha}\langle Q\rangle_{h_1+\alpha g}\\
 &=-\int_0^1d\alpha\bigl[\langle QV(g)\rangle_{h_1+\alpha g}
   -\langle Q\rangle_{h_1+\alpha g}\langle V(g)\rangle_{h_1+\alpha g}\bigr]
 \tag{1.1.12}\\
 &=-\int_0^1d\alpha\sum_{i\in\mathbb Z^2}\bigl[\langle QV(g_i)\rangle_{h_1+\alpha g}
   -\langle Q\rangle_{h_1+\alpha g}\langle V(g_i)\rangle_{h_1+\alpha g}\bigr] .
\end{align*}
We apply Theorem 1.1.7 to each term in the sum to obtain
\[
 |\langle Q\rangle_{h_2}-\langle Q\rangle_{h_1}|
 \le\sum_{i\in\mathbb Z^2\cap\operatorname{supp}g}
 O(1)e^{-m_1\operatorname{dist}(\mathcal O,\Delta_i)}
 \le O(1)e^{-m_1d/2} .
\]
As a special case, the cutoff Schwinger functions $S_h$ converge as $h\to1$.''

\emph{(b) Proof of Theorem 1.1.7, pp.~596--597.}  ``Proof of Theorem 1.1.7, assuming
Theorems 1.1.8 and 1.1.11.  Take $n=0$ in Theorem 1.1.11.  Then $Q$ is a polynomial of
degree zero, i.e., a constant.  From the choice $B=I$, it follows that
\[
 Q=\int A\,dq_h+O(1)\|A\|e^{-m_0(1-\varepsilon)T} .
\]
After rotating and translating the strip separating $Q_1$ and $Q_2$ in Theorem 1.1.7,
and after choosing $A=Q_1$, $B_T=Q_2$ in Theorem 1.1.11, the proof follows from Theorem
1.1.8, (1.1.16) and the above determination of $Q$.''  The estimate (1.1.16), p.~596,
reads ``$|\int(A-Q)B_T\,dq_h|\le O(1)\|A\|\,\|B\|e^{-\gamma T}$ uniformly in $h$.''

\emph{(c) p.~629.}  ``Theorem 1.1.11 follows from this inequality, while Theorem 1.1.8
is the special case $B=I$, $n=0$.''  The uniformity sentence of the same page is Source
Fact~\ref{sf:printed}.
\end{srcfactstar}

\emph{(viii) Scaling.} The unitary space-time dilation of \cite[p.~590]{GJS} sends
$(m_0,\lambda)\mapsto(s^{-1}m_0,s^{-2}\lambda)$; the exact reduction of the dimensionless
smallness hypothesis to the printed one is Proposition~\ref{prop:scaling}.

\emph{Scope.} Theorem 1.1.7 states that its constants are independent of $h$; the
one-sentence statement of Theorem 1.1.8 says ``uniformly as $h\uparrow1$''.  The
quantifier imported in (vi) --- one constant for the whole slab family
$\CE^{\rm slab}$ --- is the one the source states and uses, by the printed record of
item~(vii$'$), in three steps.
\emph{(a)} The source proves Theorem 1.1.1 along the affine path $h_1+\alpha(h_2-h_1)$,
$\alpha\in[0,1]$, for arbitrary cutoff functions $h_1,h_2$ of its class (p.~589:
compact support in $\mathbb R^2$, $0\le h\le1$), and applies Theorem 1.1.7 at every
cutoff of that path with one $h$-independent constant.  Our path
$h_{T,g_s}=(1-s)h_{T,g}+s\,h_{T,g'}$, $g_s=(1-s)g+sg'$ (\S\S\ref{sec:M2},
\ref{sec:G}), is this family with $h_1=h_{T,g}$, $h_2=h_{T,g'}$, $\alpha=s$.
\emph{(b)} The printed proof of Theorem 1.1.7 obtains its $h$-independent constant from
Theorem 1.1.8 and (1.1.16) \emph{at the same $h$}: the bound of Theorem 1.1.8 on
$\int Q_2\,dq_h$ multiplies the remainder $O(1)\|Q_1\|e^{-m_0(1-\varepsilon)T}$ in the
determination of $Q$.  Lemma~\ref{lem:ce2strip} and Proposition~\ref{prop:discharge}
use Theorem 1.1.7 as printed at $h=h_{T,g}\in\CE^{\rm slab}$; consistency therefore
requires the constant of Theorem 1.1.8 to be uniform on the same class, which is
\eqref{eq:UGJS}.
\emph{(c)} Source Fact~\ref{sf:printed} states the uniformity of Theorems 1.1.8 and
1.1.11 in the space cutoff $h$ explicitly, and p.~629 identifies Theorem 1.1.8 as the
case $B=I$, $n=0$ of the estimate proved there uniformly in $h$.
Hence \eqref{eq:UGJS} is imported at the scope at which the source states and uses it;
it is not a quantifier strengthened beyond the source.  It is kept as a named
hypothesis of Theorem~\ref{thm:MAIN} because its statement is assembled from four
pages of the source rather than read off one theorem, and because this manuscript does
not re-derive the cluster expansion: Appendix~\ref{app:huniform} audits the source
mechanism and proves two analytic prerequisites only.  The formerly illegible norm
prefactor was settled by page-image inspection: $K_1$ enters to the first power and the
local weight is $K_2^{2N(\Delta)}$ (gate R2,
\texttt{m4/r2-check.tex}). All constants in this manuscript use the exact reading
\eqref{eq:printednorm}.

\emph{Used in.} \S\ref{sec:reduction} (scaling), \S\ref{sec:slab}, \S\ref{sec:M2} (CE1),
\S\ref{sec:norms} (norms), \S\ref{sec:G} (CE1, CE2), Appendix~\ref{app:huniform}.

{\raggedright\emph{Dossier.} \texttt{gjs-source-05/06}, \texttt{gjs-p589-06}, \texttt{gjs-scaling-07},
\texttt{gjs-theorems-10/11/12}, \texttt{gjs-core-11/12}, \texttt{gjs-uniform-46},
\texttt{gjs-uniformity-46}, \texttt{gjs-proof-14} (E20--E26, E54).\par}
\end{assumption}

\begin{remark}[conventions lock]\label{rem:conv}
% src: m6/main.tex rem:conv (v2), transcribed.
The sharp-time restriction of the Euclidean covariance equals $C$:
$\int\frac{dk_0}{2\pi}(k_0^2+\mu(k)^2)^{-1}=(2\mu(k))^{-1}$ exactly. Hence sharp-time
$d\mu_0$-Wick ordering coincides, constant by constant, with the $C$-Wick ordering used in
$V(h)$, and sharp-time functionals of $\Phi(0,\cdot)$ are multiplication operators on $\HH$.
No Wick-mass conversion arises anywhere below.
\end{remark}

\subsection{Immediate consequences of the imports}\label{subsec:consequences}

\begin{proposition}[domain smoothing at the ground state]\label{prop:smoothing}
$\Om_g\in D(V(h))$ for every $g\in\CH$ and every real $h\in L^1\cap L^2$ of compact support.
\end{proposition}

\begin{proof}
% src: m6/main.tex prop:standard-main(i) (v2), transcribed.
By Assumption~\ref{A:domain}(i), for $t>0$ the range of $e^{-tK(g)}$ is contained in
$D(V(h))$. (Replacing $K(g)$ by $H(g)=K(g)-E(g)$ multiplies the semigroup by the nonzero
scalar $e^{tE(g)}$ and does not change its range.) Since $K(g)\Om_g=E(g)\Om_g$,
\[
  \Om_g=e^{tE(g)}e^{-tK(g)}\Om_g\in\Ran e^{-tK(g)}\subseteq D(V(h)) .\qedhere
\]
\end{proof}

\begin{proposition}[cutoff propagation, uniformly in $g$]\label{prop:propagation}
$\beta^g_t(\A(O))\subseteq\A(O_{|t|})$ for every $g\in\CH$ and $t\in\R$, with propagation
speed $\le1$ independent of $g$.
\end{proposition}

\begin{proof}
% src: m6/main.tex prop:standard-main(ii) (v2), transcribed.
The interaction automorphism $\mathrm{Ad}\,e^{itV(g)}$ has speed zero: for $A$ generated by
Weyl operators supported in $B\Subset O$ pick $\chi\in C_c^\infty(O)$, $\chi=1$ near $B$; then
$V(g)=V(\chi g)+V((1-\chi)g)$ with strongly commuting summands,
$e^{itV(\chi g)}\in\A(O)$, $e^{itV((1-\chi)g)}\in\A(B)'$, so
$\mathrm{Ad}\,e^{itV(g)}(A)\in\A(O)$; extend by strong density. The free dynamics has speed
one (Assumption~\ref{A:patch}(i)). The $n$-th Trotter approximant of $\beta^g_t(A)$ then lies
in $\A(O_{|t|})$, which is strongly closed, and the Trotter formula converges strongly
(Assumption~\ref{A:patch}(i)).
\end{proof}

\begin{lemma}[covariance and parity of the cutoff states]\label{lem:covparity}
For every $g\in\CH$ and $a\in\R$:
\begin{equation}
 \omega_{\tau_ag}=\omega_g\circ\tau_{-a} .
 \tag{T}\label{eq:T}
\end{equation}
For even $\mathscr P$: $\omega_g\circ\theta=\omega_g$.
\end{lemma}

\begin{proof}
% src: m6/main.tex prop:M0-main(ii),(iii) (v2), transcribed.
By Assumption~\ref{A:covar}, $U_aK(g)U_a^*=K(\tau_ag)$, so by simplicity
(Assumption~\ref{A:cutoff}) $\Om_{\tau_ag}=U_a\Om_g$ up to a phase, which cancels in the
state; this is \eqref{eq:T}.

For parity: $\Theta=\Gamma(-1)$ is unitary, self-adjoint, commutes with $H_0$, and for even
$\mathscr P$ satisfies $\Theta V(g)\Theta=V(g)$ (Assumption~\ref{A:covar}); hence
$\Theta\Om_g=c\Om_g$ with $c\in\{\pm1\}$ by simplicity and $\Theta^2=1$. In $Q$-space
$(\Theta\psi)(q)=\psi(-q)$, so $\Theta\Om_g>0$ a.e.\ together with $\Om_g>0$ a.e.\
(Assumption~\ref{A:cutoff}) forces $c=+1$. Thus
$\omega_g(\theta(A))=\ip{\Theta\Om_g}{A\Theta\Om_g}=\omega_g(A)$.
\end{proof}

\subsection{Standing weak-coupling assumption}\label{subsec:weak}

% src: m6/main.tex \S2.2 (v2), transcribed.
Fix source-admissible constants $M_*>0$ and $\lambda_*>0$ such that the Glimm--Jaffe--Spencer
expansion is valid at bare mass $M_*$ whenever $0\le\lambda'<\lambda_*$. Define
\[
  \varepsilon_{\mathscr P}:=\frac{\lambda_*}{M_*^2} .
\]
Throughout, $0\le\lambda/m_0^2<\varepsilon_{\mathscr P}$. By
Proposition~\ref{prop:scaling} and Remark~\ref{rem:wlog} we may, and from now on do,
assume that the pair $(m_0,\lambda)$ itself satisfies the printed source hypotheses
($m_0=M_*$ and $\lambda<\lambda_*$), so that the \emph{unit-lattice} statements of
Assumption~\ref{A:gjs} apply literally in the working coordinates; $m_1>0$ and
$C_7,C_8,K_1,K_2$ denote the printed constants. Remark~\ref{rem:wlog} records how each
conclusion transports back to the original coordinates.

% ==================================================================
\section{Dimensionless weak-coupling reduction}\label{sec:reduction}

\begin{lemma}[class compatibility]\label{lem:classes}
% [new] item, flagged per plan D8(i); elementary.
\emph{(i)} $\G\subset\CH$.
\emph{(ii)} For every $g\in\CH$ and every compact interval $I\subset\R$, the slab
cutoff $\mathbf 1_I\otimes g$ belongs to $\CE$; in particular
$h_{T,g}=\mathbf 1_{[-T,T]}\otimes g\in\CE$ and $\mathbf 1_{[0,T]}\otimes g\in\CE$.
\emph{(iii)} For $g,g'\in\CH$ and $s\in[0,1]$, the affine path $g_s:=(1-s)g+sg'$ stays in
$\CH$, and $\mathbf 1_I\otimes g_s=(1-s)\,\mathbf 1_I\otimes g+s\,\mathbf 1_I\otimes g'$
for every compact interval $I$, so that
$s\mapsto\mathcal U_{\mathbf 1_I\otimes g_s}
=(1-s)\mathcal U_{\mathbf 1_I\otimes g}+s\,\mathcal U_{\mathbf 1_I\otimes g'}$ is
affine. If moreover $g,g'\in\G$ and $\{g=1\}^\circ\cap\{g'=1\}^\circ\ne\varnothing$, then
$g_s\in\G$ with $\{g_s=1\}^\circ\supseteq\{g=1\}^\circ\cap\{g'=1\}^\circ$; in particular
$d(g_s,O)\ge\min\{d(g,O),d(g',O)\}$ for every bounded interval $O$.
\end{lemma}

\begin{proof}
(i) A function in $C_c^\infty(\R)$ with values in $[0,1]$ is measurable with compact
support. (ii) $\mathbf 1_I\otimes g$ is a product of measurable functions, takes values
in $[0,1]$, and $\supp(\mathbf 1_I\otimes g)\subseteq I\times\supp g$ is compact.
(iii) A convex combination of
$[0,1]$-valued measurable functions is $[0,1]$-valued and measurable, and
$\supp g_s\subseteq\supp g\cup\supp g'$ is compact; the identity
$\mathbf 1_I\otimes g_s=(1-s)\,\mathbf 1_I\otimes g+s\,\mathbf 1_I\otimes g'$ holds
pointwise, and $h\mapsto\mathcal U_h$ is linear in $h$ on the common $L^2(d\mu_0)$
domain guaranteed by Assumption~\ref{A:stability}. For the last statement:
$g_s\in C_c^\infty$, $0\le g_s\le1$, and
$\{g_s=1\}\supseteq\{g=1\}\cap\{g'=1\}$ (for $s\in(0,1)$ a convex combination of two
numbers in $[0,1]$ equals $1$ if and only if both do, and for $s\in\{0,1\}$ the
inclusion is trivial); passing to interiors gives the inclusion, which is a nonempty open
set by hypothesis. The distance bound follows since
$d(\cdot,O)$ is monotone in the plateau interior.
\end{proof}

\begin{proposition}[dimensionless weak-coupling reduction]\label{prop:scaling}
% src: m6/main.tex prop:scaling (v2) = lamport_euclidean_to_gap.tex <1>1, transcribed.
The standing inequality $0\le\lambda/m_0^2<\varepsilon_{\mathscr P}$ implies the source
smallness hypothesis of Assumption~\ref{A:gjs}. More precisely, the GJS unitary dilation
with $s=m_0/M_*$ sends
\[
  (m_0,\lambda)\longmapsto
  (s^{-1}m_0,s^{-2}\lambda)
  =\left(M_*,\frac{M_*^2}{m_0^2}\lambda\right),
\]
and the transformed coupling is strictly smaller than $\lambda_*$.
\end{proposition}

\begin{proof}
% src: m6/main.tex prop:scaling proof (v2), transcribed.
For $m>0$, write the Euclidean covariance, as an identity of tempered distributions, as
\[
  C_m(z)=\frac1{(2\pi)^2}\int_{\R^2}\frac{e^{ip\cdot z}}{p^2+m^2}\,dp .
\]
Since $sM_*=m_0$, the substitution $q=sp$ gives, without an omitted power of $s$,
\begin{align*}
  C_{M_*}(sz)
  &=\frac1{(2\pi)^2}\int_{\R^2}
      \frac{e^{ip\cdot sz}}{p^2+M_*^2}\,dp \\
  &=\frac1{(2\pi)^2}\int_{\R^2}
      \frac{s^{-2}e^{iq\cdot z}}{s^{-2}q^2+M_*^2}\,dq
    =C_{m_0}(z).
\end{align*}
Thus the centered generalized Gaussian fields $\Phi_{m_0}(x)$ and $\Phi_{M_*}(sx)$ have
identical joint covariances after smearing. This also fixes Wick ordering without a hidden
constant. Indeed, for every regularized centered Gaussian variable $X$,
\[
 {:}e^{tX}{:}=\exp\!\left(tX-\tfrac12t^2\E[X^2]\right),
\]
and the coefficient of $t^n$ is ${:}X^n{:}/n!$. Apply the identity to a common mollifier,
transported by $x\mapsto sx$, and then take the defining distributional limit. Equality of
the regularized covariances shows that the unitary Gaussian identification sends
${:}\mathscr P(\Phi_{m_0}(x)){:}_{m_0}$ to
${:}\mathscr P(\Phi_{M_*}(sx)){:}_{M_*}$ with the coefficients of $\mathscr P$ unchanged.
For a Euclidean cutoff $h$ the change of variables $y=sx$ therefore gives
\[
 \lambda\int_{\R^2}h(x){:}\mathscr P(\Phi_{m_0}(x)){:}_{m_0}\,dx
 =s^{-2}\lambda\int_{\R^2}h(y/s){:}\mathscr P(\Phi_{M_*}(y)){:}_{M_*}\,dy .
\]
The transformed cutoff $h_s(y):=h(y/s)$ still satisfies $0\le h_s\le1$ and has compact
support, so $h_s\in\CE$. Hence the transformed parameters are exactly those displayed in
the proposition, as stated in \cite[p.~590]{GJS}. Finally,
\[
 s^{-2}\lambda=\frac{M_*^2}{m_0^2}\lambda
 <M_*^2\varepsilon_{\mathscr P}=\lambda_* .
\]
This is precisely the source smallness condition. If the source decay rate in the scaled
coordinate $y=sx$ is $m_{1,*}>0$, then
$e^{-m_{1,*}d_y}=e^{-sm_{1,*}d_x}$; thus the rate in the original coordinates is
$m_1:=sm_{1,*}>0$. The transport of the remaining structure is
Remark~\ref{rem:wlog}.
\end{proof}

\begin{remark}[coordinate convention and transport of the conclusions]\label{rem:wlog}
% [new] item, authorized by the U4 plan amendment in strategy.md (adjudicates the U4
% finding on the lattice transport). Each claim is a one-line consequence of the
% computations in the proof of Proposition 3.2.
The dilation behind Proposition~\ref{prop:scaling} identifies the model at
$(m_0,\lambda)$ with the model at $(M_*,\lambda')$, $\lambda'=(M_*/m_0)^2\lambda
<\lambda_*$, as follows. Let $s=m_0/M_*$, $y=sx$, and let $\sigma_s$ be the scaling
automorphism of the Weyl algebra, $\sigma_s(W(f,h))=W(s^{-1}f(\cdot/s),h(\cdot/s))$
(a symplectic map). The sharp-time covariance is scale-covariant between the two masses,
$C_{M_*}(s\,z)=C_{m_0}(z)$ (the one-dimensional analogue of the substitution in the
proof of Proposition~\ref{prop:scaling}), so $\sigma_s$ carries the mass-$m_0$ Fock
state onto the mass-$M_*$ Fock state and is unitarily implemented by some $U_s$ with
$U_s\Om_0^{(m_0)}=\Om_0^{(M_*)}$; moreover $\sigma_s(\A(O))=\A(sO)$. On the
generators, $\mu_{M_*}(k/s)=s^{-1}\mu_{m_0}(k)$ (since $sM_*=m_0$) gives
$U_sH_0^{(m_0)}U_s^*=s\,H_0^{(M_*)}$, and the Euclidean Wick identification of
Proposition~\ref{prop:scaling}, restricted to sharp time (Remark~\ref{rem:conv}), gives
$U_sV_{m_0,\lambda}(g)U_s^*=s\,V_{M_*,\lambda'}(g(\cdot/s))$; hence
\[
 U_s\,K_{m_0,\lambda}(g)\,U_s^*=s\,K_{M_*,\lambda'}(g(\cdot/s)),
 \qquad g\in\CH .
\]
Consequently $g\mapsto g(\cdot/s)$ is an order isomorphism of $\G$ onto $\G$ (and a
bijection of $\CH$), cutoff ground states correspond under $U_s$, and the two cutoff
nets of states are carried onto each other. Every conclusion of
Theorem~\ref{thm:MAIN}, proved in the working (scaled) coordinates, therefore holds for
the original model with the dictionary: localization $O\leftrightarrow sO$; plateau
distance $d\leftrightarrow s\,d$, so the rate becomes $\Psi_{m_{1,*}}(s\,d)$, again
decaying faster than every inverse power of $d$; dynamics
$\alpha_t\leftrightarrow\alpha_{st}$, the light cone being form-invariant; gap
$m_1:=s\,m_{1,*}$; and the local-moment clause of item (5) transports to supports of
length $2/s$ --- equivalently, it holds for supports of length $2$ after splitting $h$
by a partition of unity into at most $\lceil s\rceil+1$ pieces of support length
$\le2/s$ (finite additivity of $V$), with constant $(\lceil s\rceil+1)$ times the
transported one. All displayed constants of \S\S\ref{sec:slab}--\ref{sec:lppl} and of
Appendices~\ref{app:constants}, \ref{app:huniform} are the constants of the working
coordinates, in which the unit-lattice geometry of Assumption~\ref{A:gjs} is literal.
\end{remark}

% ==================================================================
\section{Slab transfer}\label{sec:slab}

Throughout this section and \S\S\ref{sec:M2}, \ref{sec:G}, a cutoff $g\in\CH$ is fixed
where limits $T\to\infty$ are taken; \emph{no uniformity in $g$ beyond the printed
constants is used, and no gap is assumed}. Set
% src: m6/main.tex \S3 (v2), transcribed.
\[
\begin{gathered}
  h_{T,g}:=\mathbf 1_{[-T,T]}\otimes g\in\CE\quad(\text{Lemma \ref{lem:classes}(ii)}),\\
  \psi_s:=e^{-sH(g)}\Om_0,\qquad\hat\psi_s:=\psi_s/\|\psi_s\|,\qquad
  c_g:=\ip{\Om_g}{\Om_0}>0
\end{gathered}
\]
($c_g>0$ by Assumption~\ref{A:cutoff}: $\Om_g>0$ a.e., $\Om_0\equiv1$).

\begin{lemma}[slab FKN; bounded nonnegative arguments]\label{lem:slabFKN}
% src: m6/main.tex lem:FKN-main (v2), statement split; = lamport_euclidean_to_gap.tex <1>2<2>1.
For $T>0$, $t\in[0,2T]$, $s\in[-T,T]$, and bounded nonnegative measurable functions
$F_1,F_2$ of the time-zero field:
\begin{align}
  \int F_1\bigl(\Phi(-\tfrac t2,\cdot)\bigr)F_2\bigl(\Phi(\tfrac t2,\cdot)\bigr)dq_{h_{T,g}}
  &=\frac{\ip{F_1\psi_{T-t/2}}{e^{-tH(g)}F_2\psi_{T-t/2}}}{\|\psi_T\|^2},\label{eq:fkn2}\\
  \int F_1\bigl(\Phi(s,\cdot)\bigr)\,dq_{h_{T,g}}
  &=\frac{\ip{\psi_{T-s}}{F_1\,\psi_{T+s}}}{\|\psi_T\|^2};\label{eq:fkn1}
\end{align}
all energy shifts cancel in the ratios.
\end{lemma}

\begin{proof}
% src: m6/main.tex lem:FKN-main proof, first paragraph (v2), transcribed.
This is the Feynman--Kac--Nelson formula for $Z(h_{T,g})^{-1}e^{-\mathcal U_{h_{T,g}}}d\mu_0$
(Markov property of $d\mu_0$ along constant-time lines; Assumption~\ref{A:fkn}; the measure
is well-defined by Assumption~\ref{A:stability} and Lemma~\ref{lem:classes}(ii); the
identification of sharp-time Wick functionals with $Q$-space multiplications is
Remark~\ref{rem:conv}). The unshifted formula contains $K(g)$ and the same factor
$e^{-2TE(g)}$ in numerator and partition function; multiplying every semigroup by the
corresponding $e^{aE(g)}$ replaces $K(g)$ by $H(g)$ and cancels that factor exactly. For
bounded nonnegative $F_i$ both sides are finite.
\end{proof}

\begin{lemma}[signed bounded extension]\label{lem:slabsigned}
% src: m6/main.tex lem:FKN-main proof (v2) = lamport_euclidean_to_gap.tex <1>2<2>2.
Equations \eqref{eq:fkn2}--\eqref{eq:fkn1} hold for bounded signed measurable $F_1,F_2$.
\end{lemma}

\begin{proof}
Write $F_i=F_i^+-F_i^-$ and expand both sides into the same four (respectively two) finite
nonnegative terms of Lemma~\ref{lem:slabFKN}; bilinearity (respectively linearity) proves
the identities.
\end{proof}

\begin{lemma}[unbounded extension]\label{lem:slabunbdd}
% src: m6/main.tex lem:FKN-main proof (v2) = lamport_euclidean_to_gap.tex <1>2<2>3; audit repairs L3, L4.
\emph{(i)} For nonnegative measurable $F_1,F_2$, the identities
\eqref{eq:fkn2}--\eqref{eq:fkn1} hold in $[0,\infty]$.
\emph{(ii)} For real measurable $F$, define the clipped truncation
$F^{[\ell]}:=(-\ell)\vee(F\wedge\ell)$, $\ell>0$. Whenever $F\psi_a\in\HH$,
\[
 \|(F-F^{[\ell]})\psi_a\|^2
 =\int |F-F^{[\ell]}|^2|\psi_a|^2\,d\phi_0\longrightarrow0
 \qquad(\ell\uparrow\infty),
\]
by dominated convergence with dominator $|F|^2|\psi_a|^2$; and on path space
\[
 |F_1(\Phi(s_1,\cdot))F_2(\Phi(s_2,\cdot))|
 \le\tfrac12\bigl(F_1(\Phi(s_1,\cdot))^2+F_2(\Phi(s_2,\cdot))^2\bigr)
\]
is a dominated-convergence majorant whenever the two square moments are finite. Consequently
the signed identities used below are limits of the bounded signed identities of
Lemma~\ref{lem:slabsigned}, with no formal manipulation of products of unbounded operators.
\end{lemma}

\begin{proof}
(i) Apply Lemma~\ref{lem:slabFKN} to $F_i\wedge\ell$ and let $\ell\uparrow\infty$;
monotone convergence on both sides gives the identities in $[0,\infty]$. (ii) The
Hilbert-space limit is dominated convergence as displayed. Membership $\psi_a\in D(F)$ is
supplied at each application: finiteness of the corresponding $t=0$ identity of part (i)
over the $a$-slab, together with the slab bound \eqref{eq:UGJS}, forces $F\psi_a\in\HH$ (this
is how it is verified in Theorem~\ref{thm:M2}, Step~5, and in
Lemma~\ref{lem:transfer}); for the interaction-type multiplications it also follows
independently from Assumption~\ref{A:domain}(i). On path space,
$F_i^{[\ell]}(\Phi(s_i,\cdot))\to F_i(\Phi(s_i,\cdot))$ pointwise with the displayed
majorant, whose integrability is exactly the finiteness of the two square moments; the
working-family hypothesis \eqref{eq:UGJS} supplies those square moments at every later application
(\S\ref{sec:M2}, \S\ref{sec:G}). Dominated convergence then passes
Lemma~\ref{lem:slabsigned} to the limit on both sides.
\end{proof}

\begin{lemma}[slab vacuum; simplicity only]\label{lem:slabvac}
% src: m6/main.tex lem:vac-main (v2), transcribed; = lamport_euclidean_to_gap.tex <1>3.
$\hat\psi_s\to\Om_g$ in norm, $\|\psi_s\|\downarrow c_g$, and
$\|\psi_{s'}\|/\|\psi_s\|\in[1,1/c_g]$ for $0\le s'\le s$; the ratio tends to $1$ whenever
$s,s'\to\infty$. \emph{No isolation of the ground-state eigenvalue is used.}
\end{lemma}

\begin{proof}
% src: m6/main.tex lem:vac-main proof (v2), transcribed.
Set $P_g=|\Om_g\rangle\langle\Om_g|$. Since $P_g\Om_0=c_g\Om_g$,
\[
 \psi_s=c_g\Om_g+e^{-sH(g)}(1-P_g)\Om_0,
\]
and the two summands are orthogonal. If $\mu_g^\perp$ is the spectral measure of
$(1-P_g)\Om_0$, then
\[
 \|e^{-sH(g)}(1-P_g)\Om_0\|^2
 =\int_{(0,\infty)}e^{-2sE}\,d\mu_g^\perp(E)\longrightarrow0
\]
by dominated convergence. The measure has no atom at $0$ by simplicity alone
(Assumption~\ref{A:cutoff}); no positive distance from zero is used. It follows that
$\psi_s\to c_g\Om_g$, $\|\psi_s\|\to c_g$, and $\hat\psi_s\to\Om_g$. The spectral formula
also shows that $s\mapsto\|\psi_s\|$ is decreasing. Since
$c_g\le\|\psi_s\|\le\|\psi_0\|=1$, for $0\le s'\le s$,
\[
 1\le\frac{\|\psi_{s'}\|}{\|\psi_s\|}\le\frac1{c_g}.
\]
If both indices tend to infinity, numerator and denominator tend to $c_g>0$, so their ratio
tends to $1$.
\end{proof}

% ==================================================================
\section{The uniform local interaction moment}\label{sec:M2}

\begin{theorem}[uniform local interaction moment (M$_2$)]\label{thm:M2}
% src: lamport_euclidean_to_gap.tex <1>4; constant m4/r2-check.tex Prop. 2.1, Cor. 2.2.
Let $0\le\lambda/m_0^2<\varepsilon_{\mathscr P}$. For every $g\in\CH$ and every real
measurable $h$ with $\|h\|_\infty\le1$ and $\supp h$ contained in an interval of length
$2$,
\begin{equation}
 \|V(h)\Om_g\|\le M,\qquad
 M^2:=9\,C_8\,\mathfrak K(p')\,\lambda^2\Bigl(\sum_{n=0}^{2p'}|a_n|\Bigr)^2,
 \qquad
 \mathfrak K(p'):=K_1K_2^{8p'}(4p')! .
 \tag{M$_2$}\label{eq:M2}
\end{equation}
\end{theorem}

\begin{proof}
\emph{Step 1: the squared interaction as a GJS observable.}
% src: lamport_euclidean_to_gap.tex <1>4 (4.1), <2>1.
In Euclidean notation define
\begin{equation}
 Q_V:=\lambda^2\sum_{n,n'=0}^{2p'}a_na_{n'}
 \iint h(x)h(y)\,{:}\Phi(0,x)^n{:}\,{:}\Phi(0,y)^{n'}{:}\,dx\,dy .
 \label{eq:QV}
\end{equation}
{\sloppy $Q_V$ is exactly the multiplication function $V(h)^2$: expanding the square
of $\lambda\sum_na_n\int h(x){:}\Phi(0,x)^n{:}\,dx$ gives \eqref{eq:QV}.\par} No Wick reordering of
the product is performed or required; the observable class of Assumption~\ref{A:gjs}(ii)
permits products of individually Wick-ordered powers at equal times, and by
Remark~\ref{rem:conv} the sharp-time Wick ordering coincides with the $C$-ordering of
$V(h)$.

\emph{Step 2: cell count.}
% src: lamport_euclidean_to_gap.tex <1>4 <2>2.
The support of $h$ meets at most three half-open unit spatial cells $[j,j+1)$ in positive
measure. Indeed, if it met four distinct cells in positive measure, the extreme indices
would satisfy $j_{\max}\ge j_{\min}+3$, and the containing interval would have infimum
$<j_{\min}+1$ and supremum $>j_{\max}\ge j_{\min}+3$, hence length $>2$.

\emph{Step 3: local weight of one ordered cell pair.}
% src: lamport_euclidean_to_gap.tex <1>4 <2>3.
For one ordered pair of cells and degrees $n,n'$, the GJS local weight of
\eqref{eq:printednorm} is bounded by $K_1K_2^{8p'}(4p')!$. Indeed the total local degree is
$n+n'\le4p'$. If both variables belong to one square, the weight is
$K_1K_2^{2(n+n')}(n+n')!$; if to two squares, it is $K_1K_2^{2n}n!\,K_2^{2n'}n'!$. In both
cases
\[
 K_2^{2(n+n')}\le K_2^{8p'}\quad(K_2\ge1),\qquad
 n!\,n'!\le(n+n')!\le(4p')! .
\]

\emph{Step 4: the GJS norm of $Q_V$.}
% src: lamport_euclidean_to_gap.tex <1>4 <2>4 (4.2).
Sum over the at most $3^2=9$ ordered cell pairs and all coefficient pairs. The kernel of
the $(n,n',\Delta,\Delta')$ piece is $w(x,y)=h(x)\mathbf 1_\Delta(x)\,h(y)\mathbf
1_{\Delta'}(y)$, and $\|w\|_2=\|h\mathbf 1_\Delta\|_2\|h\mathbf 1_{\Delta'}\|_2\le1$
because $\|h\mathbf 1_\Delta\|_2^2=\int_\Delta h^2\le1$ ($\|h\|_\infty\le1$, unit cell).
Hence
\begin{equation}
 \|Q_V\|_{\rm GJS}\le
 9\,K_1K_2^{8p'}(4p')!\,\lambda^2\Bigl(\sum_{n=0}^{2p'}|a_n|\Bigr)^2=M^2/C_8 .
 \label{eq:QVnorm}
\end{equation}

\emph{Step 5: transfer at time zero.}
% src: lamport_euclidean_to_gap.tex <1>4 <2>5 (4.3); audit repair L3.
For every $T>0$, apply \eqref{eq:fkn2} at $t=0$ with the two \emph{individually
nonnegative} functions $F_1=F_2=|V(h)|$ (Lemma~\ref{lem:slabunbdd}(i)); their product is
$V(h)^2=Q_V$. Then \eqref{eq:UGJS} and \eqref{eq:QVnorm} give
\begin{equation}
 \|V(h)\hat\psi_T\|^2
 =\int V(h)^2\,dq_{h_{T,g}}\le C_8\|Q_V\|_{\rm GJS}\le M^2 .
 \label{eq:VT}
\end{equation}
In particular the left side is finite, so $\psi_T\in D(V(h))$.

\emph{Step 6: weak-graph passage to $\Om_g$.}
% src: lamport_euclidean_to_gap.tex <1>4 <2>6.
By \eqref{eq:VT}, the family $(V(h)\hat\psi_T)_{T>0}$ is bounded by $M$; choose
$T_k\to\infty$ with $V(h)\hat\psi_{T_k}\rightharpoonup y$ weakly. The graph of the closed
operator $V(h)$ is a closed linear subspace of $\HH\oplus\HH$, hence weakly closed. Since
$\hat\psi_{T_k}\to\Om_g$ strongly (Lemma~\ref{lem:slabvac}), the pair $(\Om_g,y)$ lies in
the graph; thus $\Om_g\in D(V(h))$, $y=V(h)\Om_g$, and weak lower semicontinuity of the
norm gives
\[
 \|V(h)\Om_g\|\le\liminf_k\|V(h)\hat\psi_{T_k}\|\le M .\qedhere
\]
\end{proof}

\begin{corollary}[$\varphi^4$ instance]\label{cor:M2phi4}
% src: m4/r2-check.tex Cor. 2.2, transcribed.
For $\mathscr P(\xi)=\xi^4$ ($p'=2$, $\sum_n|a_n|=1$; degrees exactly $4+4=8$),
\[
 M^2=9\cdot8!\;C_8K_1K_2^{16}\lambda^2=362\,880\;C_8K_1K_2^{16}\lambda^2 .
\]
\end{corollary}

% ==================================================================
\section{GJS norms of sharp-time polynomial observables}\label{sec:norms}

Throughout this section, $Q=\prod_{i=1}^r\varphi_0(f_i)$ with real $f_i\in C_c^\infty(\R)$,
all supported in an interval $I$ of length $L$, viewed at a common sharp time as an
observable of the class of Assumption~\ref{A:gjs}(ii). We call such $Q$ \emph{cylinder
monomials} and their linear span, together with $\mathbf 1$, \emph{cylinder polynomials}.
Set
\[
 N_I:=\#\{j\in\Z:\ |I\cap[j,j+1)|>0\} .
\]

\begin{lemma}[cell count]\label{lem:cells}
% src: lamport_euclidean_to_gap.tex <1>5 <2>1; audit repair L1.
$N_I\le\lfloor L\rfloor+2\le L+2$. The bound $L+1$ is % negative-scope mention (whitelisted): L+1
false for nonintegral $L$: the interval $I=(-0.1,0.1)$ has length $0.2$ and meets both
$[-1,0)$ and $[0,1)$ in positive measure.
\end{lemma}

\begin{proof}
Let $j_{\min}$ and $j_{\max}$ be the smallest and largest indices of cells meeting $I$ in
positive measure, so $N_I\le j_{\max}-j_{\min}+1$. Positive-measure intersection with
$[j_{\min},j_{\min}+1)$ forces $\inf I<j_{\min}+1$, and with $[j_{\max},j_{\max}+1)$ forces
$\sup I>j_{\max}$; hence
\[
 L=\sup I-\inf I>j_{\max}-(j_{\min}+1)=j_{\max}-j_{\min}-1 .
\]
The integer $j_{\max}-j_{\min}-1$ is strictly less than $L$, hence at most
$\lceil L\rceil-1\le\lfloor L\rfloor$. Consequently
$N_I\le(j_{\max}-j_{\min}-1)+2\le\lfloor L\rfloor+2$.
\end{proof}

\begin{proposition}[GJS norms of polynomial observables]\label{prop:gjsnorms}
% src: lamport_euclidean_to_gap.tex <1>5 <2>2, (5.1)-(5.2).
\begin{align}
 \|Q\|_{\rm GJS}
 &\le K_1K_2^{2r}r!\,N_I^{r/2}\prod_{i=1}^r\|f_i\|_2,\label{eq:normQ}\\
 \|Q^2\|_{\rm GJS}
 &\le K_1K_2^{4r}(2r)!\,N_I^{r}\prod_{i=1}^r\|f_i\|_2^2
   =:\mathfrak n(Q)^2 .\label{eq:normQ2}
\end{align}
\end{proposition}

\begin{proof}
Expand the kernel of $Q$ into cell assignments: decompose each
$f_i=\sum_j f_i\mathbf 1_{[j,j+1)}$, where only the $N_I$ cells meeting $I$ in positive
measure contribute. For any fixed assignment of the $r$ factors to cells, the localized
piece has, by \eqref{eq:printednorm}, the weight
$K_1\prod_\Delta K_2^{2N(\Delta)}N(\Delta)!$ with $\sum_\Delta N(\Delta)=r$; the product of
local factorials is at most the factorial of the total degree, $\prod_\Delta
N(\Delta)!\le r!$, and the $K_2$-exponents sum to $2r$. The localized kernel is the tensor
product of the restrictions, with $L^2$ norm $\prod_i\|f_i\mathbf 1_{[j_i,j_i+1)}\|_2$.
Summing over assignments ((1.1.9) of Assumption~\ref{A:gjs}(iii)) and using, for each $i$,
\[
 \sum_{j:\,|I\cap[j,j+1)|>0}\|f_i\mathbf 1_{[j,j+1)}\|_2
 \le N_I^{1/2}\|f_i\|_2
\]
by Cauchy--Schwarz over the $N_I$ cells, proves \eqref{eq:normQ}. Applying the same
argument to the $2r$ factors of $Q^2$ (degree $2r$, $K_2$-exponent $4r$, factorial
$(2r)!$, and $N_I^{1/2}$ per factor) proves \eqref{eq:normQ2}.
\end{proof}

\begin{remark}[row independence]\label{rem:row}
% src: lamport_euclidean_to_gap.tex <1>5 <2>3.
The norm of a sharp-time monomial is independent of the common sharp time $s$: at every $s$
all factors occupy one temporal row, and changing $s$ changes only that common row label,
not the spatial cell assignments, local degrees, kernel norms, or the product of weights.
This is the precise statement; arbitrary time translation does not literally preserve the
fixed lattice of unit squares.
\end{remark}

% ==================================================================
\section{The uniform gap}\label{sec:G}

Throughout this section $0\le\lambda/m_0^2<\varepsilon_{\mathscr P}$, $g\in\CH$ is
arbitrary and fixed, and $Q=\prod_{i=1}^r\varphi_0(f_i)$ is as in \S\ref{sec:norms}. Write
$Q(s)$ for the same monomial at sharp Euclidean time $s$.

\begin{lemma}[moment transfer]\label{lem:transfer}
% src: lamport_euclidean_to_gap.tex <1>6 <2>1 (6.1).
$\displaystyle\sup_{T>0}\|Q\hat\psi_T\|^2\le C_8\,\mathfrak n(Q)^2$; moreover
$\Om_g\in D(Q)$ and $\|Q\Om_g\|^2\le C_8\,\mathfrak n(Q)^2$.
\end{lemma}

\begin{proof}
Apply \eqref{eq:fkn2} at $t=0$ with $F_1=F_2=|Q|$ (Lemma~\ref{lem:slabunbdd}(i)); the
Euclidean integrand is $Q^2$ at time $0$, so \eqref{eq:UGJS}, Remark~\ref{rem:row}
and \eqref{eq:normQ2} give
\[
 \|Q\hat\psi_T\|^2=\int Q(0)^2\,dq_{h_{T,g}}\le C_8\|Q^2\|_{\rm GJS}
 \le C_8\,\mathfrak n(Q)^2
\]
for every $T$.  Choose from the bounded family $Q\widehat\psi_T$ a weakly
convergent subnet $Q\widehat\psi_{T_\alpha}\rightharpoonup y$.  Since
$\widehat\psi_{T_\alpha}\to\Om_g$ strongly and the graph of the closed
multiplication operator $Q$ is a closed linear subspace of $\HH\oplus\HH$, hence
weakly closed, $(\Om_g,y)$ lies in that graph.  Thus $\Om_g\in D(Q)$,
$y=Q\Om_g$, and weak lower semicontinuity gives
$\|Q\Om_g\|\le\liminf_\alpha\|Q\widehat\psi_{T_\alpha}\|le
C_8^{1/2}\mathfrak n(Q)$.
\end{proof}

\begin{lemma}[one-point limit]\label{lem:onepoint}
% src: lamport_euclidean_to_gap.tex <1>6 <2>2 (6.2)-(6.4); audit repair L2.
For every fixed $s\in\R$,
\[
 \int Q(s)\,dq_{h_{T,g}}\longrightarrow\ip{\Om_g}{Q\Om_g}
 \qquad(T\to\infty).
\]
\end{lemma}

\begin{proof}
For $T>|s|$, \eqref{eq:fkn1} (extended to $Q$ by Lemma~\ref{lem:slabunbdd}; the square moments
$\int Q(s)^2\,dq_{h_{T,g}}\le C_8\|Q^2\|_{\rm GJS}$ at the needed times follow from
\eqref{eq:UGJS} with Remark~\ref{rem:row}, as in Lemma~\ref{lem:transfer}) gives
\begin{equation}
 \int Q(s)\,dq_{h_{T,g}}
 =\Bigl\langle\hat\psi_{T-s},\,\frac{Q\psi_{T+s}}{\|\psi_T\|}\Bigr\rangle
   \frac{\|\psi_{T-s}\|}{\|\psi_T\|}.
 \label{eq:onepoint}
\end{equation}
The first vector tends strongly to $\Om_g$ and the last factor tends to $1$
(Lemma~\ref{lem:slabvac}). Write the middle vector as
\begin{equation}
 \frac{Q\psi_{T+s}}{\|\psi_T\|}
 =\frac{\|\psi_{T+s}\|}{\|\psi_T\|}\,Q\hat\psi_{T+s} .
 \label{eq:middle}
\end{equation}
If $s\ge0$, the scalar ratio in \eqref{eq:middle} is at most $1$; if $s<0$, it is at most
$c_g^{-1}$, both by Lemma~\ref{lem:slabvac}. (The inequality
$\|\psi_{T+s}\|\le\|\psi_T\|$ holds only for $s\ge0$; the two signs must be treated
separately.) Thus \eqref{eq:middle} is bounded in either case, by
Lemma~\ref{lem:transfer}. Every weak subsequential limit of \eqref{eq:middle} equals
$Q\Om_g$, by weak closedness of the graph of $Q$, strong convergence
$\hat\psi_{T+s}\to\Om_g$, and because the scalar ratio tends to $1$. Hence the whole middle
vector converges weakly to $Q\Om_g$, and \eqref{eq:onepoint} converges to
$\ip{\Om_g}{Q\Om_g}$.
\end{proof}

\begin{lemma}[two-point lower semicontinuity]\label{lem:twopoint}
% src: lamport_euclidean_to_gap.tex <1>6 <2>3 (6.5)-(6.6).
For every $t\ge0$, with the Euclidean integrals over slabs of half-width $T\ge t/2$
(so that \eqref{eq:fkn2} applies),
\[
 \ip{Q\Om_g}{e^{-tH(g)}Q\Om_g}
 \le\liminf_{T\to\infty}\int Q(-t/2)\,Q(t/2)\,dq_{h_{T,g}} .
\]
\end{lemma}

\begin{proof}
With $u_T:=Q\psi_{T-t/2}/\|\psi_T\|$, the extension of \eqref{eq:fkn2} to $Q$
(Lemma~\ref{lem:slabunbdd}; square moments at the two times by
\eqref{eq:UGJS} and Remark~\ref{rem:row}, as in Lemma~\ref{lem:transfer}) is the
exact equality
\[
 \int Q(-t/2)\,Q(t/2)\,dq_{h_{T,g}}=\bigl\|e^{-tH(g)/2}u_T\bigr\|^2 .
\]
The ratio bound of Lemma~\ref{lem:slabvac} and Lemma~\ref{lem:transfer} make $(u_T)_T$
bounded.  Indeed,
\[
 u_T=\frac{\|\psi_{T-t/2}\|}{\|\psi_T\|}\,
       Q\widehat\psi_{T-t/2},
\]
the scalar factor tends to $1$ and is bounded by $c_g^{-1}$, and
$\widehat\psi_{T-t/2}\to\Om_g$ strongly.  From any subnet of $(u_T)$ choose a
weakly convergent further subnet, say $u_{T_\alpha}\rightharpoonup u$.  Along a
further subnet the bounded vectors $Q\widehat\psi_{T_\alpha-t/2}$ converge weakly
to the same $u$.  Since the graph of the closed linear operator $Q$ is a closed
linear subspace of $\HH\oplus\HH$, it is weakly closed; hence
$(\Om_g,u)$ belongs to that graph and $u=Q\Om_g$.  Every weak cluster point is
therefore $Q\Om_g$, so the full bounded net satisfies
$u_T\rightharpoonup Q\Om_g$. Boundedness of $e^{-tH(g)/2}$ preserves the weak convergence
$e^{-tH(g)/2}u_T\rightharpoonup e^{-tH(g)/2}Q\Om_g$, and weak lower semicontinuity of the
norm gives the claim.
\end{proof}

\begin{lemma}[CE2 across a horizontal strip]\label{lem:ce2strip}
% src: lamport_euclidean_to_gap.tex <1>6 <2>4 (6.7).
For $t>2$ and every $h\in\CE$,
\[
 \Bigl|\int Q(-t/2)Q(t/2)\,dq_h
 -\int Q(-t/2)\,dq_h\int Q(t/2)\,dq_h\Bigr|
 \le C_7e^{2m_1}e^{-m_1t}\,\|Q\|_{\rm GJS}^2 .
\]
\end{lemma}

\begin{proof}
Use the finite localization decomposition prescribed by GJS (1.1.9):
\[
 Q=\sum_{a\in A}Q_a,
 \qquad
 \|Q\|_{\rm GJS}=\sum_{a\in A}\|Q_a\|_{\rm GJS},
\]
where every $Q_a$ is a localized monomial and $A$ is finite because the kernels
of the cylinder polynomial have compact spatial support and finite degree.  By
linearity of expectation,
\[
 \operatorname{Cov}_{dq_h}(Q(-t/2),Q(t/2))
 =\sum_{a,b\in A}
   \operatorname{Cov}_{dq_h}(Q_a(-t/2),Q_b(t/2)).
\]
The localization squares of $Q_a(-t/2)$ lie in the row of unit squares meeting
$x^0=-t/2$, and those of $Q_b(t/2)$ lie in the row meeting $x^0=t/2$.
The open strip $\{|x^0|<t/2-1\}$ has width $t-2$ and separates every such pair.
Source Fact~\ref{sf:CE2}, applied separately to each $(a,b)$, and the triangle
inequality therefore give
\begin{align*}
 \left|\operatorname{Cov}_{dq_h}(Q(-t/2),Q(t/2))\right|
 &\le C_7e^{-m_1(t-2)}
   \sum_{a,b\in A}\|Q_a(-t/2)\|_{\rm GJS}
                    \|Q_b(t/2)\|_{\rm GJS}\\
 &=C_7e^{-m_1(t-2)}
   \left(\sum_a\|Q_a\|_{\rm GJS}\right)^2\\
 &=C_7e^{2m_1}e^{-m_1t}\|Q\|_{\rm GJS}^2,
\end{align*}
where the second equality uses the row-independence proved in
Remark~\ref{rem:row}.  No application of the localized CE2 theorem to an
unlocalized sum is made.
\end{proof}

\begin{proposition}[semigroup decay of centered observables]\label{prop:semidecay}
% src: lamport_euclidean_to_gap.tex <1>6 <2>5 (6.8).
Let $\psi_Q:=Q\Om_g-\ip{\Om_g}{Q\Om_g}\Om_g$. Then
\[
 \ip{\psi_Q}{e^{-tH(g)}\psi_Q}\le\widehat C_Q\,e^{-m_1t}\qquad(t\ge0),
 \qquad
 \widehat C_Q:=\max\bigl\{C_7e^{2m_1}\|Q\|_{\rm GJS}^2,\;e^{3m_1}\|\psi_Q\|^2\bigr\} .
\]
\end{proposition}

\begin{proof}
Let $t>2$. Since $e^{-tH(g)}\Om_g=\Om_g$ and
$\ip{\Om_g}{Q\Om_g}$ is real,
\[
 \ip{\psi_Q}{e^{-tH(g)}\psi_Q}
 =\ip{Q\Om_g}{e^{-tH(g)}Q\Om_g}-\ip{\Om_g}{Q\Om_g}^2 .
\]
Write $\int Q(-t/2)Q(t/2)\,dq_{h_{T,g}}=a_T+b_T$, where
$b_T=\int Q(-t/2)\,dq_{h_{T,g}}\int Q(t/2)\,dq_{h_{T,g}}$ and $a_T$ is the connected
difference. Then, by Lemma~\ref{lem:twopoint} and
$\liminf_T(a_T+b_T)\le\sup_T|a_T|+\lim_Tb_T$,
\[
 \ip{Q\Om_g}{e^{-tH(g)}Q\Om_g}
 \le\sup_T|a_T|+\lim_Tb_T
 \le C_7e^{2m_1}e^{-m_1t}\|Q\|_{\rm GJS}^2+\ip{\Om_g}{Q\Om_g}^2 ,
\]
using Lemma~\ref{lem:ce2strip} for $\sup_T|a_T|$ and Lemma~\ref{lem:onepoint} at
$s=\mp t/2$ fixed for $\lim_Tb_T$. Subtracting the square proves the bound for $t>2$. For
$0\le t\le3$, positivity of $H(g)$ gives
$\ip{\psi_Q}{e^{-tH(g)}\psi_Q}\le\|\psi_Q\|^2\le e^{3m_1}\|\psi_Q\|^2e^{-m_1t}$. The two
cases cover $[0,\infty)$ and both bounds are $\le\widehat C_Qe^{-m_1t}$.
\end{proof}

\begin{lemma}[exponential moments]\label{lem:expmom}
% src: lamport_euclidean_to_gap.tex <1>7 <2>1 (7.1)-(7.2); audit repair L5.
Let $d\mu_g:=|\Om_g|^2\,d\phi_0$ and $X_f:=\varphi_0(f)$ for real
$f\in C_c^\infty(\R)$, $f\ne0$, supported in an interval of length $L$, and put
$K':=K_2^{2}N_I^{1/2}\|f\|_2>0$. Then
\[
 \E_{\mu_g}\bigl[X_f^{2n}\bigr]\le C_8K_1\,(K')^{2n}(2n)! \qquad(n\ge1),
\]
and consequently
\[
 \E_{\mu_g}e^{s|X_f|}
 \le2\max\{1,C_8K_1\}\sum_{n=0}^\infty(sK')^{2n}<\infty
 \qquad(0\le s<1/K') .
\]
The factor $C_8K_1$ affects only the prefactor, not the radius $1/K'$.
\end{lemma}

\begin{proof}
$\E_{\mu_g}[X_f^{2n}]=\|X_f^n\Om_g\|^2\le C_8\,\mathfrak n(X_f^n)^2$ by
Lemma~\ref{lem:transfer} applied to $Q=X_f^n$; and by \eqref{eq:normQ2} with $r=n$ and all
$f_i=f$,
\[
 \mathfrak n(X_f^n)^2=K_1K_2^{4n}(2n)!\,N_I^{n}\|f\|_2^{2n}
 =K_1(K')^{2n}(2n)! .
\]
For the exponential bound: $e^{s|x|}\le e^{sx}+e^{-sx}=2\cosh(sx)$, and
\[
 \E_{\mu_g}\cosh(sX_f)=\sum_{n\ge0}\frac{s^{2n}}{(2n)!}\E_{\mu_g}[X_f^{2n}]
 \le\max\{1,C_8K_1\}\sum_{n\ge0}(sK')^{2n},
\]
where the maximum covers the $n=0$ term $\E_{\mu_g}[1]=1$. The series converges for
$sK'<1$.
\end{proof}

\begin{theorem}[density of polynomial vectors]\label{thm:dense}
% src: lamport_euclidean_to_gap.tex <1>7 <2>2-<2>5.
Cylinder polynomials are dense in $L^2(d\mu_g)$. Equivalently, since $\Om_g>0$ a.e.\ makes
multiplication by $\Om_g$ a unitary from $L^2(d\mu_g)$ onto $\HH$, the linear span of
$\Om_g$ and the vectors $Q\Om_g$ ($Q$ a cylinder monomial) is dense in $\HH$.
\end{theorem}

\begin{proof}
\emph{Step 1: vanishing conditional expectations.} Let $\zeta\in L^2(d\mu_g)$ be orthogonal
to all cylinder polynomials. Fix real $e_1,\dots,e_k\in C_c^\infty(\R)$ and
define
\[
 F(z_1,\dots,z_k):=
 \E_{\mu_g}\Bigl[\overline\zeta\,\exp\Bigl(\sum_{j=1}^kz_jX_{e_j}\Bigr)\Bigr] .
\]
By Cauchy--Schwarz, H\"older, and Lemma~\ref{lem:expmom}, $F$ is defined and analytic in a
nonempty polystrip $\{|\mathrm{Re}\,z_j|<s_0\}$ about $i\R^k$; differentiation under the
integral is justified by the exponential dominator of Lemma~\ref{lem:expmom}. Every
derivative of $F$ at zero is an inner product of $\zeta$ with a cylinder polynomial, hence
zero. Repeated application of the one-variable identity theorem gives $F=0$ throughout the
connected polystrip, in particular on $i\R^k$. Thus the finite complex measure
\[
 B\longmapsto\E_{\mu_g}\bigl[\overline\zeta\,
 \mathbf 1_{\{(X_{e_1},\dots,X_{e_k})\in B\}}\bigr]
\]
on $\R^k$ has vanishing Fourier transform, so it is zero by uniqueness of Fourier
transforms of finite measures; equivalently
$\E_{\mu_g}[\zeta\mid\sigma(X_{e_1},\dots,X_{e_k})]=0$.

\emph{Step 2: exhaustion.} Choose $(e_j)_{j\ge1}\subseteq C_c^\infty(\R)$, real and
dense in real $L^2(\R)$ ($C_c^\infty$ is dense there). The time-zero
covariance $C=(2\mu)^{-1}$ is bounded, so
\[
 \E_{\phi_0}|X_f-X_e|^2=\ip{f-e}{C(f-e)}\le\|C\|\,\|f-e\|_2^2 .
\]
For an approximating sequence $e_{j_l}\to f$ in $L^2$, a subsequence of $X_{e_{j_l}}$
converges $\phi_0$-almost surely, hence $\mu_g$-almost surely because
$\mu_g\ll\phi_0$. Therefore every $X_f$ is measurable with respect to
$\bigvee_k\sigma(X_{e_1},\dots,X_{e_k})$ modulo null sets, and this sigma-algebra exhausts
the cylinder sigma-algebra.

\emph{Step 3: martingale convergence.} The conditional expectations of Step~1 vanish for
every $k$ and converge in $L^2(d\mu_g)$ to the conditional expectation on the exhausted
sigma-algebra, which is $\zeta$ itself. Hence $\zeta=0$, proving density.
\end{proof}

\begin{theorem}[uniform gap (G)]\label{thm:G}
% src: lamport_euclidean_to_gap.tex <1>8; audit repair L6.
Let $0\le\lambda/m_0^2<\varepsilon_{\mathscr P}$. For every $g\in\CH$,
\begin{equation}
 \spec H(g)\subseteq\{0\}\cup[m_1,\infty),
 \tag{G}\label{eq:G}
\end{equation}
and $0$ is a simple eigenvalue.
\end{theorem}

\begin{proof}
Fix $0\le\beta<m_1$ and let $P_\beta:=\mathbf 1_{[0,\beta]}(H(g))$.

\emph{Step 1: $P_\beta\psi_Q=0$ for every centered polynomial vector.} By the spectral
theorem and Proposition~\ref{prop:semidecay}, for every $t\ge0$,
\[
 e^{-\beta t}\|P_\beta\psi_Q\|^2
 \le\ip{\psi_Q}{e^{-tH(g)}\psi_Q}
 \le\widehat C_Q\,e^{-m_1t} .
\]
Multiply by $e^{\beta t}$ and let $t\to\infty$; since $m_1-\beta>0$,
$\|P_\beta\psi_Q\|=0$. (The restriction $\beta\ge0$ is necessary: for $\beta<0$, $P_{[0,\beta]}=0$, and
Step~2 below fails.)

\emph{Step 2: $P_\beta\Om_g=\Om_g$}, since $H(g)\Om_g=0$ and $0\in[0,\beta]$.

\emph{Step 3: $P_\beta=|\Om_g\rangle\langle\Om_g|$.} By Theorem~\ref{thm:dense}, the
complex span of $\Om_g$ and the centered vectors $\psi_Q$ is dense in $\HH$. Steps~1--2
show that $P_\beta$ maps this dense subspace into $\C\Om_g$; continuity of the bounded
projection extends the containment to $\HH$, and $P_\beta\Om_g=\Om_g$ shows the range is
exactly $\C\Om_g$.

\emph{Step 4: no spectrum in $(0,m_1)$.} If an open interval $J\Subset(0,m_1)$ carried a
nonzero spectral projection, choose $\beta<m_1$ above the right endpoint of $J$; then
$\mathbf 1_J(H(g))\le P_\beta-|\Om_g\rangle\langle\Om_g|=0$, a contradiction. Hence
\eqref{eq:G} holds; simplicity of $0$ is Assumption~\ref{A:cutoff}.
\end{proof}

% ==================================================================
\section{Affine-path regularity}\label{sec:path}

Throughout \S\S\ref{sec:path}--\ref{sec:lppl}, $0\le\lambda/m_0^2<\varepsilon_{\mathscr P}$,
and $\gamma>0$ denotes a constant for which
\begin{equation}
 \spec H(g)\subseteq\{0\}\cup[\gamma,\infty),\qquad \dim\ker H(g)=1,
 \qquad\text{for every }g\in\CH ;
 \tag{G$_\gamma$}\label{eq:Ggamma}
\end{equation}
by Theorem~\ref{thm:G} this holds with $\gamma=m_1$, and $\gamma$ is instantiated there in
\S\ref{sec:limit}. Fix $g,g'\in\G$ equal to $1$ on a common nonempty open interval, and put
% src: lamport_gap_to_main.tex <1>1, setting.
\[
 g_s:=(1-s)g+sg',\quad v:=g'-g,\quad K_s:=K(g_s),\quad E_s:=E(g_s),\quad H_s:=K_s-E_s,
\]
\[
 \Om_s:=\Om_{g_s},\qquad P_s:=|\Om_s\rangle\langle\Om_s|,\qquad Q_s:=1-P_s .
\]
By Lemma~\ref{lem:classes}(iii), $g_s\in\G\subset\CH$ for every $s\in[0,1]$, so all results
of \S\S\ref{sec:slab}--\ref{sec:G} apply to every $K_s$.

\begin{proposition}[path regularity of the semigroup]\label{prop:pathreg}
% src: lamport_gap_to_main.tex <1>1 <2>1.
For each $T>0$, the family $s\mapsto e^{-TK_s}$ is norm-continuous on $[0,1]$ and
norm-$C^1$ on $(0,1)$, with derivative given by the FKN kernel of $-\,\mathcal U_vw_s$
in the notation of the proof.
\end{proposition}

\begin{proof}
On free Euclidean path space over the slab $[0,T]$ define
% src: lamport_gap_to_main.tex <1>1 <2>1, transcribed.
\[
 \mathcal U_s:=\lambda\int_0^T\!\!\int_\R g_s(x)\,{:}\mathscr P(\Phi(t,x)){:}\,dx\,dt,
 \qquad w_s:=e^{-\mathcal U_s},\qquad
 \mathcal U_v:=\mathcal U_1-\mathcal U_0 ,
\]
so that $s\mapsto\mathcal U_s=(1-s)\mathcal U_0+s\,\mathcal U_1$ is affine
(Lemma~\ref{lem:classes}(iii)). For bounded time-zero $f,h$, the interacting FKN formula of
Assumption~\ref{A:fkn} reads
\begin{equation}
 \ip{f}{e^{-TK_s}h}=\int\overline{(J_Tf)}\,(J_0h)\,w_s\,d\mu_0 ,
 \label{eq:fknpath}
\end{equation}
with $J_t$ the sharp-time embeddings.

\emph{Step 1: the pairing bound.} Set $a:=e^{-m_0T}$ and
\[
 q_T:=\frac2{1-a},\quad r_T:=\frac2{1+a},\quad p_T:=1+a,\quad p_T':=1+a^{-1} .
\]
Then $1/r_T+1/q_T=1$, $r_Tp_T=2$, and $(p_T'-1)/(p_T-1)=a^{-2}=e^{2m_0T}$, so
Assumption~\ref{A:hyper} gives $\|\mathsf P_TF\|_{p_T'}\le\|F\|_{p_T}$ for the free Markov
semigroup $\mathsf P_T$ on $L^2(d\phi_0)$. The Markov property of $d\mu_0$ along constant
time lines and H\"older's inequality (note $1/p_T+1/p_T'=1$) yield
\begin{align*}
 \|(J_Tf)(J_0h)\|_{r_T}^{r_T}
 &=\int |h|^{r_T}\,\mathsf P_T\bigl(|f|^{r_T}\bigr)\,d\phi_0
 \le\bigl\||h|^{r_T}\bigr\|_{p_T}\bigl\|\mathsf P_T(|f|^{r_T})\bigr\|_{p_T'}
 \le\|h\|_2^{r_T}\,\|f\|_2^{r_T},
\end{align*}
using $\||h|^{r_T}\|_{p_T}=\|h\|_{r_Tp_T}^{r_T}=\|h\|_2^{r_T}$ and hypercontractivity on
the second factor. Hence
\begin{equation}
 \|(J_Tf)(J_0h)\|_{r_T}\le\|f\|_2\|h\|_2 .
 \label{eq:pairing}
\end{equation}

\emph{Step 2: $L^a$ control and continuity of $s\mapsto w_s$.} Finite-volume stability
(Assumption~\ref{A:stability}, applicable since the slab cutoffs lie in $\CE$ by
Lemma~\ref{lem:classes}(ii)) gives $\mathcal U_v\in L^p(d\mu_0)$ for every finite $p$ and
$e^{-\varrho\,\mathcal U_i}\in L^1(d\mu_0)$ for every finite $\varrho>0$, $i=0,1$. Since
$\mathcal U_s$ is affine, H\"older's inequality with exponents
$\frac1{1-s},\frac1s$ gives
\begin{equation}
 \|w_s\|_\varrho^\varrho\le
 \Bigl(\int e^{-\varrho\,\mathcal U_0}d\mu_0\Bigr)^{1-s}
 \Bigl(\int e^{-\varrho\,\mathcal U_1}d\mu_0\Bigr)^{s} ,
 \label{eq:wsbound}
\end{equation}
a bound uniform in $s\in[0,1]$ for each $\varrho$. If $s_n\to s$, then
$\mathcal U_{s_n}\to\mathcal U_s$ in probability (indeed in $L^2$), so
$w_{s_n}\to w_s$ in probability; together with \eqref{eq:wsbound} at a larger exponent,
Vitali's convergence theorem gives $w_{s_n}\to w_s$ in $L^\varrho$ for every finite
$\varrho$.

\emph{Step 3: the difference quotient.} The calculus identity
\begin{equation}
 \frac{w_{s+\delta}-w_s}{\delta}
 =-\,\mathcal U_v\int_0^1w_{s+\vartheta\delta}\,d\vartheta
 \label{eq:diffquot}
\end{equation}
holds pointwise on path space. Choose finite $p,\varrho$ with
$q_T^{-1}=p^{-1}+\varrho^{-1}$. By H\"older, the $L^{q_T}$ distance of
\eqref{eq:diffquot} from $-\,\mathcal U_vw_s$ is at most
$\|\mathcal U_v\|_p\int_0^1\|w_{s+\vartheta\delta}-w_s\|_\varrho\,d\vartheta$, which tends
to $0$ as $\delta\to0$ for $0<s<1$ by Step~2 (dominated convergence in $\vartheta$, using
\eqref{eq:wsbound}). Thus $s\mapsto w_s$ is $L^{q_T}$-differentiable on $(0,1)$ with
derivative $-\,\mathcal U_vw_s$, and $L^{q_T}$-continuous on $[0,1]$.

\emph{Step 4: duality.} By \eqref{eq:fknpath}, \eqref{eq:pairing}, and H\"older with the
conjugate pair $(r_T,q_T)$,
\[
 \bigl|\ip{f}{(e^{-TK_s}-e^{-TK_r})h}\bigr|
 \le\|(J_Tf)(J_0h)\|_{r_T}\,\|w_s-w_r\|_{q_T}
 \le\|f\|_2\|h\|_2\,\|w_s-w_r\|_{q_T},
\]
so $\|e^{-TK_s}-e^{-TK_r}\|\le\|w_s-w_r\|_{q_T}$, and the same estimate applied to the
difference quotients transfers the $L^{q_T}$ regularity of Step~3 to the operator norm.
This proves the asserted norm continuity and differentiability.
\end{proof}

\begin{proposition}[projection and eigenvector regularity; derivative formulas]\label{prop:projreg}
% src: lamport_gap_to_main.tex <1>1 <2>2-<2>3.
\emph{(i)} $s\mapsto P_s$ is norm-continuous on $[0,1]$ and norm-$C^1$ on $(0,1)$.
\emph{(ii)} Locally on $(0,1)$ one may choose normalized $C^1$ vectors $\Om_s$ in the
parallel gauge $\ip{\Om_s}{\dot\Om_s}=0$, and then
\begin{equation}
 \dot E_s=\ip{\Om_s}{V(v)\Om_s},\qquad
 \dot\Om_s=-H_s^{-1}Q_sV(v)\Om_s ,
 \label{eq:derivform}
\end{equation}
with $\|H_s^{-1}Q_s\|\le\gamma^{-1}$.
\end{proposition}

\begin{proof}
\emph{(i)} Let $L_s:=e^{-TK_s}$ for a fixed $T>0$. By Assumption~\ref{A:cutoff}, $L_s$ has
the simple largest eigenvalue $\lambda_s=e^{-TE_s}=\|L_s\|$, and by \eqref{eq:Ggamma},
\[
 \spec L_s\setminus\{\lambda_s\}\subseteq[0,e^{-T\gamma}\lambda_s] .
\]
The relative spectral separation $1-e^{-T\gamma}$ is uniform in $s$. Since $s\mapsto L_s$
is norm-continuous, so is $s\mapsto\|L_s\|=\lambda_s$; a Riesz contour of relative radius
independent of $s$ around $\lambda_s$ therefore represents $P_s$ as a norm-convergent
Cauchy integral of the resolvent of $L_s$, locally in $s$, and the norm regularity of
$L_s$ (Proposition~\ref{prop:pathreg}) passes to $P_s$ through the contour integral. This
gives exactly the stated regularity of $P_s$.

\emph{(ii)} Let $\mathscr C$ be the common core of Assumption~\ref{A:domain}(ii). For
$\xi\in\mathscr C$, $K_s\xi=K_{s_0}\xi+(s-s_0)V(v)\xi$. Differentiating the weak
eigenvalue identity $\ip{K_s\xi}{\Om_s}=E_s\ip{\xi}{\Om_s}$ at $s=s_0$ (all terms
differentiable: $\Om_s$ is $C^1$ by (i) and the gauge; and
$E_s=-T^{-1}\log\lambda_s$ is $C^1$ because
$\lambda_s=\ip{\Om_s}{L_s\Om_s}>0$ is a composition of the $C^1$ maps of (i) and
Proposition~\ref{prop:pathreg}) gives
\begin{equation}
 \ip{H_s\xi}{\dot\Om_s}=\ip{\xi}{(\dot E_s-V(v))\Om_s} ,
 \label{eq:weakeigen}
\end{equation}
where $\Om_s\in D(V(v))$ by Proposition~\ref{prop:smoothing}, so the right side is
well-defined. Since $\mathscr C$ is a core for the self-adjoint $H_s$,
\eqref{eq:weakeigen} says $\dot\Om_s\in D(H_s)$ and
$H_s\dot\Om_s=(\dot E_s-V(v))\Om_s$. Pairing with $\Om_s$ and using
$H_s\Om_s=0$, $\ip{\Om_s}{\dot\Om_s}=0$ proves the first formula in
\eqref{eq:derivform}. The parallel gauge puts $\dot\Om_s\in Q_s\HH$, and \eqref{eq:Ggamma}
makes $H_s^{-1}Q_s$ bounded by $\gamma^{-1}$; applying $H_s^{-1}Q_s$ to
$H_s\dot\Om_s=(\dot E_s-V(v))\Om_s$ and using $Q_s\Om_s=0$ proves the second formula.
\end{proof}

% ==================================================================
\section{Filter, spectral separation, and clustering}\label{sec:filter}

\begin{lemma}[exact inverse-energy filter]\label{lem:filter}
% src: lamport_gap_to_main.tex <1>2.
Choose an even $\chi\in C^\infty(\R;[0,1])$ with $\chi=0$ on $[-\gamma/2,\gamma/2]$ and
$\chi=1$ outside $(-\gamma,\gamma)$, and set
\begin{equation}
 \widehat W_\gamma(E):=\begin{cases}\chi(E)/E,&E\ne0,\\0,&E=0,\end{cases}
 \qquad
 W_\gamma(t):=\frac1{2\pi}\int_\R e^{-itE}\,\widehat W_\gamma(E)\,dE .
 \label{eq:filterdef}
\end{equation}
Then:
\emph{(i)} $W_\gamma\in L^1(\R)$ and, for every $q\ge0$,
$\int_\R(1+|t|)^q|W_\gamma(t)|\,dt<\infty$; in particular
\[
 T_\gamma(r):=\int_{|t|>r}|W_\gamma(t)|\,dt
 \le C_{q,\gamma}(1+r)^{-q}\qquad(q\ge0,\ r\ge0).
\]
\emph{(ii)} If $H$ is self-adjoint with
$\spec H\subseteq\{0\}\cup[\gamma,\infty)$, then
\begin{equation}
 H^{-1}(1-P_{\ker H})=\int_\R W_\gamma(t)\,e^{itH}(1-P_{\ker H})\,dt ,
 \label{eq:filterid}
\end{equation}
a norm-convergent Bochner integral.
\end{lemma}

\begin{proof}
% src: lamport_gap_to_main.tex <1>2 <2>2-<2>3, transcribed.
Throughout, $\widehat W_\gamma\in L^2(\R)$ (it vanishes on
$[-\gamma/2,\gamma/2]$, is bounded by $2/\gamma$, and decays like $|E|^{-1}$), and
\eqref{eq:filterdef} is understood as its $L^2$-inverse Fourier transform; the symmetric
improper integrals below converge (Dirichlet's test) and represent it almost everywhere,
which justifies the manipulations.

\emph{(i)} Oddness of $\widehat W_\gamma$ gives
\[
 W_\gamma(t)=-\frac{i}{\pi}\int_0^\infty\frac{\sin(tE)\,\chi(E)}E\,dE .
\]
For $0<|t|\le1$: split at $E=\gamma$. The transition part
$\int_{\gamma/2}^{\gamma}|\sin(tE)|\chi(E)E^{-1}dE$ is absolutely bounded by a constant
depending only on $\gamma$; for the tail $\int_\gamma^\infty\sin(tE)E^{-1}dE$, substitute
$u=|t|E$ and use the uniform boundedness of $\int_c^d u^{-1}\sin u\,du$ over
$0<c<d$ (Dirichlet's test). Hence $W_\gamma$ is bounded near $t=0$, so integrable there.
For $|t|\ge1$: every derivative $\widehat W_\gamma^{(k)}$ with $k\ge1$ is integrable on
$\R$ (it is smooth, supported off $(-\gamma/2,\gamma/2)$, and decays like $|E|^{-1-k}$),
so $k$ integrations by parts in \eqref{eq:filterdef} give
\[
 |W_\gamma(t)|\le\frac{\|\widehat W_\gamma^{(k)}\|_{L^1}}{2\pi|t|^{k}} .
\]
Choosing $k>q+1$ proves the weighted integrability and, by Markov's inequality applied to
the weight $(1+|t|)^q$, the tail bound for $T_\gamma$.

\emph{(ii)} By construction $\widehat W_\gamma(E)=E^{-1}$ on $[\gamma,\infty)$ and
$\widehat W_\gamma(0)=0$. Fourier inversion and the spectral theorem give
\eqref{eq:filterid} on the spectral subspace of $[\gamma,\infty)$ and on $\ker H$
separately; the integral converges in norm because $W_\gamma\in L^1$ and
$\|e^{itH}(1-P_{\ker H})\|\le1$.
\end{proof}

\begin{lemma}[spectral separation]\label{lem:separation}
% src: lamport_gap_to_main.tex <1>3; variation bound renamed M_0 -> \mathfrak m per the freeze.
Let
\[
 F(t)=\int_{[\gamma,\infty)}e^{itE}\,d\rho(E),\qquad
 G(t)=\int_{(-\infty,-\gamma]}e^{itE}\,d\sigma(E),
\]
with complex measures obeying $|\rho|(\R),|\sigma|(\R)\le\mathfrak m$, and suppose
$F=G$ on $(-R,R)$ with $\gamma R\ge1$. Then
\begin{equation}
 |F(0)|\le2\,\mathfrak m\,e^{-\gamma R/2} .
 \label{eq:sep}
\end{equation}
\end{lemma}

\begin{proof}
\emph{Step 1: the Gaussian separator.} For $a,\varepsilon>0$ define
\[
 f_{a,\varepsilon}(t):=\frac{1}{2\pi i}\,\frac{e^{-t^2/(2a)}}{t-i\varepsilon} .
\]
The transform of $(2\pi i)^{-1}(t-i\varepsilon)^{-1}$ is
$\mathbf 1_{(0,\infty)}(E)\,e^{-\varepsilon E}$ (one-pole contour integral);
multiplication by the Gaussian in $t$ convolves this with the normalized Gaussian, so
\begin{equation}
 \check f_{a,\varepsilon}(E)
 =\int_0^\infty\sqrt{\frac a{2\pi}}\,e^{-a(E-E')^2/2}\,e^{-\varepsilon E'}\,dE'
 \in[0,1],
 \label{eq:septransform}
\end{equation}
which tends to the standard normal distribution function $\Phi(\sqrt aE)$ as
$\varepsilon\downarrow0$.

\emph{Step 2: the tail of the separator.} For $t\ge R$, $|t-i\varepsilon|\ge t$ and
$t^{-1}\le t/R^2$, so
\begin{equation}
 \int_{|t|\ge R}|f_{a,\varepsilon}(t)|\,dt
 \le\frac1{\pi R^2}\int_R^\infty te^{-t^2/(2a)}\,dt
 =\frac{a}{\pi R^2}\,e^{-R^2/(2a)} .
 \label{eq:septail}
\end{equation}

\emph{Step 3: exchange of test function.} Pair $f_{a,\varepsilon}$ with $F-G$. By
hypothesis the pairing is supported in $|t|\ge R$; Fubini (justified by the finite total
variations and $f_{a,\varepsilon}\in L^1$) and then $\varepsilon\downarrow0$ with
dominated convergence in energy give
\begin{equation}
 \Bigl|\int\Phi(\sqrt aE)\,d\rho(E)-\int\Phi(\sqrt aE)\,d\sigma(E)\Bigr|
 \le\frac{2\,\mathfrak m\,a}{\pi R^2}\,e^{-R^2/(2a)} .
 \label{eq:sepexchange}
\end{equation}

\emph{Step 4: optimization.} On the support of $\rho$, $E\ge\gamma$, so
$1-\Phi(\sqrt aE)\le\tfrac12e^{-a\gamma^2/2}$; on the support of $\sigma$, $E\le-\gamma$,
so $\Phi(\sqrt aE)\le\tfrac12e^{-a\gamma^2/2}$. Inserting and subtracting
$\int d\rho=F(0)$ in \eqref{eq:sepexchange},
\[
 |F(0)|\le \mathfrak m\,e^{-a\gamma^2/2}
 +\frac{2\,\mathfrak m\,a}{\pi R^2}\,e^{-R^2/(2a)} .
\]
Set $a=R/\gamma$: both exponentials become $e^{-\gamma R/2}$, and
$2a/(\pi R^2)=2/(\pi\gamma R)\le2/\pi<1$ by $\gamma R\ge1$; the two terms sum to at most
$2\,\mathfrak m\,e^{-\gamma R/2}$, proving \eqref{eq:sep}.
\end{proof}

\begin{lemma}[clustering with one unbounded local argument]\label{lem:cluster}
% src: lamport_gap_to_main.tex <1>4.
Let $g\in\CH$, write $H=H(g)$, $\Om=\Om_g$, $\beta_t=\beta^g_t$, and assume
\eqref{eq:Ggamma}. Let $B\subset\R$ be compact, let $X=X^*$ be affiliated with $\A(U)$ for
every bounded open $U\supseteq B$, with $\Om\in D(X)$, and let $A\in\A(O)$ for a bounded
interval $O$. If $R:=\dist(B,O)>0$ satisfies $\gamma R\ge1$, then
\begin{equation}
 \bigl|\ip{X\Om}{A\Om}-\ip{X\Om}{\Om}\ip{\Om}{A\Om}\bigr|
 \le2\,\|X\Om\|\,\|A\|\,e^{-\gamma R/2} .
 \label{eq:cluster}
\end{equation}
\end{lemma}

\begin{proof}
Put $P_0=|\Om\rangle\langle\Om|$ and define
\[
 F(t):=\ip{X\Om}{(e^{itH}-P_0)A\Om},\qquad
 G(t):=\ip{A^*\Om}{(e^{-itH}-P_0)X\Om} .
\]

\emph{Step 1: spectral supports and variations.} Inserting the spectral resolution of $H$
on $(1-P_0)\HH$ and using \eqref{eq:Ggamma}, $F$ and $G$ are Fourier transforms of complex
measures supported in $[\gamma,\infty)$ and $(-\infty,-\gamma]$ respectively, with total
variations at most $\mathfrak m:=\|X\Om\|\,\|A\|$ (Cauchy--Schwarz on the polar
decomposition of the spectral measures).

\emph{Step 2: locality identity.} For $|t|<R$: by
Proposition~\ref{prop:propagation}, $\beta_t(A)\in\A(O_{|t|})$, and
$\dist(B,O_{|t|})\ge R-|t|>0$, so there is a bounded open $U\supseteq B$ disjoint from
$O_{|t|}$. Affiliation of $X$ with $\A(U)$ implies that
$\beta_t(A)\in\A(O_{|t|})\subseteq\A(U)'$ preserves $D(X)$ and commutes with $X$ on it.
Since $e^{-itH}\Om=\Om$,
\[
 \ip{X\Om}{e^{itH}A\Om}
 =\ip{\Om}{X\beta_t(A)\Om}
 =\ip{\Om}{\beta_t(A)X\Om}
 =\ip{A^*\Om}{e^{-itH}X\Om} .
\]
The $P_0$ terms agree because $\ip{\Om}{X\Om}$ is real ($X=X^*$, $\Om\in D(X)$). Hence
$F=G$ on $(-R,R)$.

\emph{Step 3: conclusion.} Lemma~\ref{lem:separation} gives
$|F(0)|\le2\mathfrak m\,e^{-\gamma R/2}$, and
$F(0)=\ip{X\Om}{A\Om}-\ip{X\Om}{\Om}\ip{\Om}{A\Om}$.
\end{proof}

% ==================================================================
\section{The one-path LPPL estimate}\label{sec:lppl}

The setting of \S\ref{sec:path} remains in force: $g,g'\in\G$ equal to $1$ on a common
nonempty open interval, $g_s=(1-s)g+sg'$, $v=g'-g$; note $|v|\le1$.

\begin{lemma}[state derivative]\label{lem:statederiv}
% src: lamport_gap_to_main.tex <1>5 <2>1 (3.15); filter sign +1/E.
For self-adjoint $A\in B(\HH)$ and $0<s<1$,
\begin{equation}
 \frac d{ds}\,\omega_{g_s}(A)
 =-2\,\mathrm{Re}\int_\R W_\gamma(t)
 \Bigl[\ip{V(v)\Om_s}{\beta^{g_s}_t(A)\Om_s}
 -\ip{V(v)\Om_s}{\Om_s}\,\omega_{g_s}(A)\Bigr]dt ,
 \label{eq:statederiv}
\end{equation}
and the integral converges absolutely.
\end{lemma}

\begin{proof}
Since $A$ is bounded and self-adjoint, Proposition~\ref{prop:projreg}(ii) gives
\[
 \frac d{ds}\,\omega_{g_s}(A)
 =2\,\mathrm{Re}\ip{\dot\Om_s}{A\Om_s}
 =-2\,\mathrm{Re}\ip{V(v)\Om_s}{H_s^{-1}Q_sA\Om_s},
\]
using \eqref{eq:derivform} and self-adjointness of $H_s^{-1}Q_s$. Insert the exact filter
\eqref{eq:filterid} for $H=H_s$ (admissible by \eqref{eq:Ggamma}), applied to the vector
$Q_sA\Om_s$: since $e^{itH_s}\Om_s=\Om_s$,
\[
 e^{itH_s}Q_sA\Om_s
 =e^{itH_s}\bigl(A\Om_s-\omega_{g_s}(A)\Om_s\bigr)
 =\beta^{g_s}_t(A)\Om_s-\omega_{g_s}(A)\Om_s ,
\]
which yields \eqref{eq:statederiv}. The integral is absolutely convergent because
$W_\gamma\in L^1$ (Lemma~\ref{lem:filter}) and the bracket is bounded by
$2\|A\|\,\|V(v)\Om_s\|$; here $V(v)\Om_s$ exists by Proposition~\ref{prop:smoothing}
($v$ has compact support).
\end{proof}

\begin{lemma}[localization of the perturbation]\label{lem:localize}
% src: lamport_gap_to_main.tex <1>5 <2>2 (3.16).
Choose a smooth partition of unity $\sum_{j\in\Z}\eta_j=1$ with $0\le\eta_j\le1$ and
$\supp\eta_j\subset(j-1,j+1)$, and set $h_j:=\eta_jv$, $V_j:=V(h_j)$, $B_j:=\supp h_j$.
Then only finitely many $h_j$ are nonzero,
\begin{equation}
 V(v)\xi=\sum_jV_j\xi
 \quad\text{on common domains},\qquad
 \|V_j\Om_s\|\le M\quad(0\le s\le1).
 \label{eq:localize}
\end{equation}
\end{lemma}

\begin{proof}
Finitely many $h_j$ are nonzero because $v$ has compact support; the finite additivity is
the convention of \S\ref{subsec:model}. Each $h_j$ is real and measurable with
$\|h_j\|_\infty\le\|\eta_j\|_\infty\|v\|_\infty\le1$ and
$\supp h_j\subseteq[j-1,j+1]$, an interval of length $2$. Since $g_s\in\CH$
(Lemma~\ref{lem:classes}) and $0\le\lambda/m_0^2<\varepsilon_{\mathscr P}$,
Theorem~\ref{thm:M2} applies to the cutoff $g_s$ and gives $\|V_j\Om_s\|\le M$.
\end{proof}

\begin{lemma}[two-regime bound]\label{lem:tworegime}
% src: lamport_gap_to_main.tex <1>6 <2>2 (3.17).
Let $A=A^*\in\A(O)$ with $\|A\|\le1$, and for a nonzero $h_j$ put
$r_j:=\dist(B_j,O)$ and
\[
 C_{j,s}(t):=\ip{V_j\Om_s}{\beta^{g_s}_t(A)\Om_s}
 -\ip{V_j\Om_s}{\Om_s}\,\omega_{g_s}(A) .
\]
If $\gamma r_j/2\ge1$, then
\begin{equation}
 |C_{j,s}(t)|\le
 \begin{cases}
 2Me^{-\gamma r_j/4},&|t|\le r_j/2,\\[2pt]
 2M,&t\in\R .
 \end{cases}
 \label{eq:tworegime}
\end{equation}
\end{lemma}

\begin{proof}
For arbitrary $t$, bound each of the two terms by
$\|V_j\Om_s\|\,\|A\|\le M$ (Lemma~\ref{lem:localize}), giving $2M$. For $|t|\le r_j/2$:
Proposition~\ref{prop:propagation} places $\beta^{g_s}_t(A)$ in $\A(O_{|t|})$, and
$\dist(B_j,O_{|t|})\ge r_j-|t|\ge r_j/2$, with $\gamma\cdot\dist(B_j,O_{|t|})\ge
\gamma r_j/2\ge1$. Apply Lemma~\ref{lem:cluster} with $H=H(g_s)$, $X=V_j$ (affiliated with
every $\A(U)$, $U\supseteq B_j$, by \S\ref{subsec:model}; $\Om_s\in D(V_j)$ by
Proposition~\ref{prop:smoothing}), and the observable $\beta^{g_s}_t(A)\in\A(O_{|t|})$;
note $\ip{\Om_s}{\beta^{g_s}_t(A)\Om_s}=\omega_{g_s}(A)$ by invariance of $\Om_s$. This
gives
\[
 |C_{j,s}(t)|
 \le2\|V_j\Om_s\|\,\|\beta^{g_s}_t(A)\|\,
 e^{-\gamma\,\dist(B_j,O_{|t|})/2}
 \le2Me^{-\gamma r_j/4} .\qedhere
\]
\end{proof}

\begin{proposition}[one-path estimate]\label{prop:onepath}
% src: lamport_gap_to_main.tex <1>6 (3.18)-(3.19).
Let $d\ge d_0:=4+2/\gamma$ and suppose $g\equiv g'\equiv1$ on $O_d$. Then
\begin{equation}
 \|(\omega_g-\omega_{g'})|_{\A(O)}\|
 \le32M\sum_{n\ge\lfloor d\rfloor}
 \Bigl(\|W_\gamma\|_1e^{-\gamma n/4}+T_\gamma(n/2)\Bigr) .
 \label{eq:onepath}
\end{equation}
\end{proposition}

\begin{proof}
\emph{Step 1: all perturbation cells are far from $O$.} On $O_d$, $v=g'-g=0$, and $O_d$ is
open, so $\supp v\subseteq\R\setminus O_d$; hence every nonzero $h_j$ has
$r_j=\dist(B_j,O)\ge d\ge d_0$. In particular
$\gamma r_j/2\ge\gamma d_0/2=2\gamma+1\ge1$, so Lemma~\ref{lem:tworegime} applies to every
$j$.

\emph{Step 2: derivative bound.} Fix $A=A^*\in\A(O)$, $\|A\|\le1$. By
Lemmas~\ref{lem:statederiv} and \ref{lem:localize},
\[
 \frac d{ds}\,\omega_{g_s}(A)
 =-2\,\mathrm{Re}\sum_j\int_\R W_\gamma(t)\,C_{j,s}(t)\,dt .
\]
Split each $t$-integral at $|t|=r_j/2$ and use \eqref{eq:tworegime}: the inner region
contributes at most $\|W_\gamma\|_1\cdot2Me^{-\gamma r_j/4}$, the outer at most
$T_\gamma(r_j/2)\cdot2M$. With the overall factor $2$,
\begin{equation}
 \Bigl|\frac d{ds}\,\omega_{g_s}(A)\Bigr|
 \le4M\sum_j\Bigl(\|W_\gamma\|_1e^{-\gamma r_j/4}+T_\gamma(r_j/2)\Bigr) .
 \label{eq:derivbound}
\end{equation}

\emph{Step 3: shell count.} At most eight indices $j$ have $r_j\in[n,n+1)$, for any
integer $n\ge\lfloor d\rfloor$. Indeed, write $O=(a,b)$. Since $B_j\subseteq[j-1,j+1]$ has
length $\le2$, avoids $O_d$, and $O_d$ is an interval of length $>2d\ge8$, the set $B_j$
lies wholly on one side of $O_d$. On the right side, compactness of $B_j$ gives
$x_j\in B_j$ with $x_j-b=r_j$; then $|x_j-j|\le1$ and $r_j\in[n,n+1)$ force
$j\in[b+n-1,b+n+2]$, an interval containing at most four integers. The left side
contributes at most four more.

\emph{Step 4: integrate and sum shells.} Both functions of $r_j$ in
\eqref{eq:derivbound} are decreasing in $r_j$, so grouping the indices into the shells of
Step~3 gives
\[
 \Bigl|\frac d{ds}\,\omega_{g_s}(A)\Bigr|
 \le4M\cdot8\sum_{n\ge\lfloor d\rfloor}
 \Bigl(\|W_\gamma\|_1e^{-\gamma n/4}+T_\gamma(n/2)\Bigr) .
\]
Integrate over $s\in[\varepsilon,1-\varepsilon]$, where differentiability holds
(Proposition~\ref{prop:projreg}), and let $\varepsilon\downarrow0$ using the endpoint norm
continuity of $s\mapsto P_s$; this bounds $|\omega_g(A)-\omega_{g'}(A)|$ by the right side
of \eqref{eq:onepath}. Finally, $\omega_g-\omega_{g'}$ is a hermitian functional, so its
norm on $\A(O)$ is the supremum of $|\omega_g(A)-\omega_{g'}(A)|$ over self-adjoint
contractions $A\in\A(O)$.
\end{proof}

\begin{lemma}[the rate function is super-polynomial]\label{lem:rate}
% src: lamport_gap_to_main.tex <1>7 (3.20)-(3.21).
Define
\begin{equation}
 \Psi_\gamma(d):=
 \begin{cases}
 \displaystyle
 32M\sum_{n\ge\lfloor d\rfloor}
 \Bigl(\|W_\gamma\|_1e^{-\gamma n/4}+T_\gamma(n/2)\Bigr), & d\ge d_0,\\[6pt]
 \max\Bigl\{2,\ \text{the same expression}\Bigr\}, & 0\le d<d_0 .
 \end{cases}
 \label{eq:Psidef}
\end{equation}
Then for every $N\ge1$ there is $C_{N,\gamma,M}<\infty$ with
\begin{equation}
 \Psi_\gamma(d)\le C_{N,\gamma,M}\,(1+d)^{-N}\qquad(d\ge0) .
 \label{eq:Psitail}
\end{equation}
In particular $\|(\omega_g-\omega_{g'})|_{\A(O)}\|\le\Psi_\gamma(d)$ whenever
$g\equiv g'\equiv1$ on $O_d$, $d\ge0$.
\end{lemma}

\begin{proof}
For $d\ge d_0$: the exponential part sums to
$\|W_\gamma\|_1e^{-\gamma\lfloor d\rfloor/4}(1-e^{-\gamma/4})^{-1}
\le C_\gamma'e^{-\gamma d/4}$, which is $O_N((1+d)^{-N})$ for every $N$. For the second
part choose $q=N+2$ in the tail bound of Lemma~\ref{lem:filter}(i):
$T_\gamma(n/2)\le C_{q,\gamma}(1+n/2)^{-q}\le C_{q,\gamma}2^{q}(1+n)^{-q}$, and
\[
 \sum_{n\ge\lfloor d\rfloor}(1+n)^{-N-2}
 \le\int_{\lfloor d\rfloor-1}^\infty(1+x)^{-N-2}\,dx
 =\frac{\lfloor d\rfloor^{-N-1}}{N+1}
 \le\frac{C_N''}{(1+d)^{N+1}} ,
\]
using $\lfloor d\rfloor\ge d-1\ge\tfrac35(1+d)$ for $d\ge d_0\ge4$. For $0\le d<d_0$:
$\Psi_\gamma$ is nonincreasing there, so the left side of \eqref{eq:Psitail} is bounded
by $\Psi_\gamma(0)=\max\{2,\,32M\sum_{n\ge0}(\|W_\gamma\|_1e^{-\gamma n/4}
+T_\gamma(n/2))\}<\infty$, while the right side is bounded below by
$C_{N,\gamma,M}(1+d_0)^{-N}$; enlarging $C_{N,\gamma,M}$ covers this range. The final assertion combines Proposition~\ref{prop:onepath} with the
trivial bound $\|\omega_g-\omega_{g'}\|\le2$ for $d<d_0$.
\end{proof}

% ==================================================================
\section{Convergence of the full net; local normality}\label{sec:limit}

From here on the abstract gap parameter is instantiated: by Theorem~\ref{thm:G},
\eqref{eq:Ggamma} holds with $\gamma=m_1$, and we write
$\Psi_{m_1}:=\Psi_\gamma|_{\gamma=m_1}$ with $M$ the constant of Theorem~\ref{thm:M2}.

\begin{theorem}[full-net convergence with rate]\label{thm:cauchy}
% src: lamport_gap_to_main.tex <1>8.
Let $0\le\lambda/m_0^2<\varepsilon_{\mathscr P}$. There is a state $\omega_\infty$ on $\A$
such that, for every bounded open interval $O$,
\[
 \lim_{g\in\G}\bigl\|(\omega_g-\omega_\infty)|_{\A(O)}\bigr\|=0 ,
\]
and
\begin{equation}
 \|(\omega_g-\omega_\infty)|_{\A(O)}\|\le\Psi_{m_1}(d)
 \qquad\text{whenever }g\equiv1\text{ on }O_d .
 \label{eq:rate}
\end{equation}
\end{theorem}

\begin{proof}
\emph{Step 1: the net is locally Cauchy.} Fix $O$ and $\varepsilon>0$. By
Lemma~\ref{lem:rate}, choose $d\ge d_0$ with $\Psi_{m_1}(d)<\varepsilon$, and choose
$g_0\in\G$ whose plateau interior contains $O_d$. If $g,g'\succcurlyeq g_0$, then
$\{g=1\}^\circ\supseteq\{g_0=1\}^\circ\supseteq O_d$ and likewise for $g'$; in particular
$g\equiv g'\equiv1$ on $O_d$ and their common plateau interior is nonempty, so
Proposition~\ref{prop:onepath} (with $\gamma=m_1$) applies and gives
\begin{equation}
 \|(\omega_g-\omega_{g'})|_{\A(O)}\|\le\Psi_{m_1}(d)<\varepsilon .
 \label{eq:cauchyeps}
\end{equation}

\emph{Step 2: limit functionals and compatibility.} Each dual $\A(O)^*$ is complete, so
\eqref{eq:cauchyeps} produces a norm limit $\varphi_O$ of $(\omega_g|_{\A(O)})_{g\in\G}$.
If $O_1\subseteq O_2$, then restriction to $\A(O_1)$ of the $\A(O_2)$-limit is again the
norm limit of $\omega_g|_{\A(O_1)}$, hence equals $\varphi_{O_1}$ by uniqueness of norm
limits. The compatible family $(\varphi_O)_O$ defines a linear functional $\omega_\infty$
on the norm-dense union $\bigcup_O\A(O)$, with $|\omega_\infty(A)|\le\|A\|$,
$\omega_\infty(1)=1$, and $\omega_\infty(A^*A)=\lim_g\omega_g(A^*A)\ge0$ for local $A$; it
extends uniquely by norm continuity to $\A$, and the extension is a state (a norm-one
positive unital functional on a norm-dense subalgebra extends to a state).

\emph{Step 3: the rate.} Fix $g\equiv1$ on $O_d$. The subset
$\{g'\in\G:g'\equiv1\text{ on }O_d\}$ is cofinal in $\G$ (every plateau can be enlarged to
contain $O_d$). For every such $g'$, Proposition~\ref{prop:onepath} gives
$\|(\omega_g-\omega_{g'})|_{\A(O)}\|\le\Psi_{m_1}(d)$ (via Lemma~\ref{lem:rate} for
$d<d_0$); letting $g'$ run through this cofinal subset, $\omega_{g'}\to\omega_\infty$ in
the norm of $\A(O)^*$, and the bound passes to the limit, proving \eqref{eq:rate}.
\end{proof}

\begin{proposition}[local normality]\label{prop:normal}
% src: lamport_gap_to_main.tex <1>9 (3.23).
$\omega_\infty|_{\A(O)}$ is normal for every bounded open interval $O$.
\end{proposition}

\begin{proof}
Let $0\le A_i\uparrow A$ be a bounded increasing net in $\A(O)$ with supremum $A$. Given
$\delta>0$, choose $g\in\G$ with
$\|(\omega_g-\omega_\infty)|_{\A(O)}\|<\delta$ (Theorem~\ref{thm:cauchy}). Then
\[
 0\le\omega_\infty(A)-\omega_\infty(A_i)
 =\bigl[\omega_\infty-\omega_g\bigr](A-A_i)+\omega_g(A-A_i)
 \le2\delta\|A\|+\omega_g(A-A_i),
\]
using $\|A-A_i\|\le2\|A\|$. The cutoff state $\omega_g$ is a vector state, hence normal on
the von Neumann algebra $\A(O)$, so $\omega_g(A-A_i)\to0$ along the net. Therefore
$\limsup_i[\omega_\infty(A)-\omega_\infty(A_i)]\le2\delta\|A\|$; let
$\delta\downarrow0$. Thus $\omega_\infty(A_i)\uparrow\omega_\infty(A)$, which is normality
on $\A(O)$.
\end{proof}

% ==================================================================
\section{Ground state, invariances, and the Main Theorem}\label{sec:ground}

\begin{proposition}[exact spectral passage]\label{prop:spectral}
% src: m6/main.tex prop:M0-main(i) (v2), transcribed; = lamport_gap_to_main.tex <1>10.
If $(g_n)\subseteq\G$ satisfies $d(g_n,O)\to\infty$ for every bounded interval $O$, and
$\omega_{g_n}\to\omega$ in norm on every $\A(O)$, then $\omega$ is a ground state of
$(\A,\alpha)$ in the spectral sense: $\omega\circ\alpha_t=\omega$ for all $t$; the
function $t\mapsto\omega(A\alpha_t(B))$ is continuous and bounded for all $A,B\in\A$; and
\[
 \int_\R f(t)\,\omega(A\alpha_t(B))\,dt=0
\]
for every $f\in L^1(\R)$ whose transform $\check f(E)=\int f(t)e^{itE}\,dt$ vanishes on
$[0,\infty)$.
\end{proposition}

\begin{proof}
% src: m6/main.tex prop:M0-main proof (i) (v2), transcribed.
\emph{Exact identity.} Fix a bounded interval $O$, $A,B\in\A(O)$,
$\mathfrak m:=\|A\|\|B\|$, and $g\in\G$; let
$\rho_g(\cdot):=\ip{A^*\Om_g}{P_{H(g)}(\cdot)B\Om_g}$, a complex measure with
$\supp\rho_g\subseteq\spec H(g)\subseteq[0,\infty)$ and $|\rho_g|(\R)\le\mathfrak m$, and
$G_g(t):=\int e^{itE}\,d\rho_g(E)$. For $|t|<d(g,O)$, the cutoff equals one on
$O_{|t|}$, so by \eqref{eq:P} and $e^{-itH(g)}\Om_g=\Om_g$,
\[
  \omega_g\bigl(A\alpha_t(B)\bigr)
  =\ip{\Om_g}{Ae^{itH(g)}Be^{-itH(g)}\Om_g}
  =\ip{A^*\Om_g}{e^{itH(g)}B\Om_g}=G_g(t) .
\]
Moreover $\int_\R f\,G_g\,dt=\int_{[0,\infty)}\check f\,d\rho_g=0$ for the stated class of
$f$ (Fubini, justified by $|\rho_g|(\R)\le\mathfrak m$ and $f\in L^1$; $\check f=0$ on the
support).

\emph{Invariance.} For $A\in\A(O)$ and fixed $t$, take $n$ so large that
$d(g_n,O)>|t|$. Then \eqref{eq:P} and $e^{-itH(g_n)}\Om_{g_n}=\Om_{g_n}$ give directly
\[
 \omega_{g_n}(\alpha_t(A))
 =\ip{\Om_{g_n}}{e^{itH(g_n)}Ae^{-itH(g_n)}\Om_{g_n}}
 =\omega_{g_n}(A).
\]
Both $\alpha_t(A)$ and $A$ belong to $\A(O_{|t|})$, so norm convergence on
$\A(O_{|t|})$ passes the identity to the limit:
$\omega(\alpha_t(A))=\omega(A)$, and by density on $\A$.

\emph{Continuity and vanishing.} Fix $A,B\in\A(O)$ and $T>0$. For $n$ with
$d(g_n,O)>T$ and $|t|\le T$: $A\alpha_t(B)\in\A(O_T)$ with norm $\le\mathfrak m$, so
\[
  \sup_{|t|\le T}\bigl|\omega(A\alpha_t(B))-\omega_{g_n}(A\alpha_t(B))\bigr|
  \le\mathfrak m\,\bigl\|(\omega-\omega_{g_n})|_{\A(O_T)}\bigr\|\longrightarrow0 :
\]
$t\mapsto\omega(A\alpha_t(B))$ is a locally uniform limit of the continuous functions
$G_{g_n}$, hence continuous, bounded by $\mathfrak m$. For $f$ as stated and
$\varepsilon(T):=\int_{|t|>T}|f|$:
\[
  \int_\R f\,\omega(A\alpha_\cdot(B))
  =\underbrace{\int_{|t|\le T}f\,[\omega-\omega_{g_n}](A\alpha_t(B))}_{\textstyle
  \le\|f\|_1\mathfrak m\|(\omega-\omega_{g_n})|_{\A(O_T)}\|}
  +\underbrace{\int_{|t|\le T}f\,G_{g_n}}_{\textstyle=\,-\int_{|t|>T}f\,G_{g_n}}
  +\int_{|t|>T}f\,\omega(A\alpha_t(B)) ,
\]
where the middle identity uses $\int_\R f\,G_{g_n}\,dt=0$. Hence the modulus is at most
$\|f\|_1\mathfrak m\|(\omega-\omega_{g_n})|_{\A(O_T)}\|+2\mathfrak m\,\varepsilon(T)$;
let $n\to\infty$, then $T\to\infty$. Joint norm continuity in $(A,B)$ extends the
conclusion from local $A,B$ to $\A$.
\end{proof}

The Main Theorem now follows by assembly.

\begin{theorem}[Main Theorem]\label{thm:MAIN}
% src: m6/main.tex thm:MAIN (v2), transcribed; import references adapted to (A1)-(A8);
% item (5) quantifies over C_H where m6 stated G -- a strengthening delivered by
% thm:G/thm:M2 (recorded in part_i/README.md, G6).
Assume $U_{\rm GJS}$, namely Assumption~\ref{A:gjs}(vi), and let
$0\le\lambda/m_0^2<\varepsilon_{\mathscr P}$. Then there is a state $\omega_\infty$ on $\A$
such that:
\begin{enumerate}
\item[\rm(1)] \emph{(full-net convergence with rate)} for every bounded open interval $O$,
      $\lim_{g\in\G}\|(\omega_g-\omega_\infty)|_{\A(O)}\|=0$, and
      $\|(\omega_g-\omega_\infty)|_{\A(O)}\|\le\Psi_{m_1}(d)$ whenever $g\equiv1$ on $O_d$,
      with $\Psi_{m_1}$ decaying faster than every inverse power;
\item[\rm(2)] \emph{(local normality)} $\omega_\infty|_{\A(O)}$ is normal for every $O$;
\item[\rm(3)] \emph{(ground state)} $\omega_\infty$ is a ground state of $(\A,\alpha)$ in the
      spectral sense;
\item[\rm(4)] \emph{(invariances)} $\omega_\infty\circ\tau_a=\omega_\infty$ for all $a\in\R$;
      and for even $\mathscr P$, $\omega_\infty\circ\theta=\omega_\infty$;
\item[\rm(5)] \emph{(spectra and local moment)}
      \[
       \spec H(g)\subseteq\{0\}\cup[m_1,\infty)\quad(g\in\CH),
       \qquad \dim\ker H(g)=1,
      \]
      and $\sup_{g\in\CH}\|V(h)\Om_g\|\le M$ for real $h$ with $\|h\|_\infty\le1$ supported
      in an interval of length $2$.
\end{enumerate}
The nonstandard weak-coupling estimates imported are printed GJS Theorem 1.1.7 and
the moment estimate of Theorem 1.1.8 in its working-family form $U_{\rm GJS}$
(Assumption~\ref{A:gjs}(vi), imported at the scope at which the source states and uses
it; Scope paragraph and item~(vii$'$) there).  Appendix~\ref{app:huniform} supplies
supporting scope evidence and analytic prerequisites but is not an independent proof
of the cluster expansion. In addition, the proof uses the
separately identified standard cutoff-construction inputs of
Assumptions~\ref{A:cutoff}--\ref{A:covar}. Proposition~\ref{prop:scaling} verifies that the
dimensionless weak-coupling hypothesis implies the source smallness condition. The theorem is conditional
on exactly this list, including $U_{\rm GJS}$; no other external input is used.
\end{theorem}

\begin{proof}
% src: m6/main.tex thm:MAIN proof (v2), transcribed; refs adapted.
(5) is Theorem~\ref{thm:G} and Theorem~\ref{thm:M2}. These supply \eqref{eq:Ggamma} with
$\gamma=m_1$ and \eqref{eq:M2}, so Theorem~\ref{thm:cauchy} and
Proposition~\ref{prop:normal} give (1)--(2). For (3): choose $\widetilde g_n\in\G$ with
$\widetilde g_n\equiv1$ on $(-n,n)$; then $d(\widetilde g_n,O)\to\infty$ for every $O$
and, by (1), $\omega_{\widetilde g_n}\to\omega_\infty$ in norm on every $\A(O)$;
Proposition~\ref{prop:spectral} applies and gives (3) together with the
$\alpha$-invariance. For (4): translation is an order isomorphism of $\G$ with cofinal
range, so by \eqref{eq:T} and (1),
\[
 \omega_\infty(A)=\lim_g\omega_{\tau_ag}(A)
 =\lim_g\omega_g(\tau_{-a}A)
 =\omega_\infty(\tau_{-a}A)
\]
for local $A$, and by density on $\A$. The $\Z_2$ statement is the local norm limit of
Lemma~\ref{lem:covparity}: $\omega_g(\theta(A))=\omega_g(A)$ for $A\in\A(O)$, and
$\theta(A)\in\A(O)$ ($\theta$ maps $W(f,h)$ to $W(-f,-h)$ and hence preserves each local
algebra), so (1) passes the identity to $\omega_\infty$, and density extends it to $\A$.
\end{proof}

\begin{corollary}[the Glimm--Jaffe vacuum is canonical at weak coupling]\label{cor:canonical}
% src: m6/main.tex cor:canonical (v2), transcribed.
Every $w^*$ limit point of the (averaged or unaveraged) finite-volume ground states
constructed in \cite[\S2]{GJ3} coincides on each $\A(O)$ with $\omega_\infty$; in
particular the construction of \cite{GJ3} is subsequence- and averaging-independent in the
weak-coupling region.
\end{corollary}

\begin{proof}
For the unaveraged net this is Theorem~\ref{thm:MAIN}(1). For the averaged states of
\cite[(2.3)--(2.5)]{GJ3}: fix $g\in\G$ with $g\equiv1$ on $[-3,3]$ and set
$g_n:=g(\cdot/n)\in\G$; let $\sigma_{a'}$ denote the space-translation automorphisms
used there (our $\tau_{\epsilon a'}$, with $\epsilon\in\{\pm1\}$ the direction
convention), and let $\bar h\ge0$, $\supp\bar h\subseteq[-1,1]$, $\int\bar h=1$, be the
averaging weight. Then
\begin{align*}
  \omega_n(A)
  &=\frac1n\int\ip{\Om_{g_n}}{\sigma_{a'}(A)\Om_{g_n}}\,
      \bar h(a'/n)\,da'
  =\int_{-1}^{1}\bar h(u)\,\omega_{g_n}\bigl(\tau_{\epsilon nu}A\bigr)\,du
  =\int_{-1}^{1}\bar h(u)\,\omega_{\tau_{-\epsilon nu}g_n}(A)\,du ,
\end{align*}
by \eqref{eq:T} (the estimate below is invariant under $\epsilon\to-\epsilon$). Since
$g_n\equiv1$ on $[-3n,3n]$, the shifted cutoff $\tau_{-\epsilon nu}g_n\in\G$ equals $1$
on an interval containing $[-2n,2n]$ for every $|u|\le1$; hence for $O\subseteq[-R,R]$,
$d(\tau_{-\epsilon nu}g_n,O)\ge2n-R$ \emph{uniformly in $u$}, and
Theorem~\ref{thm:MAIN}(1) gives
\[
  \|(\omega_n-\omega_\infty)|_{\A(O)}\|
  \le\int_{-1}^1\bar h(u)\,\bigl\|(\omega_{\tau_{-\epsilon nu}g_n}-\omega_\infty)|_{\A(O)}\bigr\|\,du
  \le\Psi_{m_1}(2n-R)\longrightarrow0 .
\]
No translation invariance of $\omega_\infty$ is needed for this bound.
\end{proof}

% ==================================================================
\section{The two-phase dichotomy}\label{sec:dichotomy}

The uniform gap of Theorem~\ref{thm:G} is incompatible with a fully classified two-phase
structure. The incompatibility is stated as a conditional theorem: \emph{all four}
hypotheses below are load-bearing, and (B4) in particular cannot be omitted --- dropping
it was the one proof-breaking defect (P1) found by the independent audit of the
predecessor manuscript (Appendix~\ref{app:provenance}).

\begin{hypothesis}[structure of a broken phase]\label{hyp:broken}
% src: m6/main.tex cor:nobreak items (B1)-(B4) (v2); m3/m3.tex hyp:broken states (B4)
% two-sidedly (segment identification) -- only the one-sided inclusion used here.
$\mathscr P$ is even, and:
\begin{enumerate}
\item[\rm(B1)] there are translation-invariant ground states $\omega_+$ and $\omega_-$ of
      $(\A,\alpha)$ such that, for all bounded local $X,Y$,
      \[
       \omega_\pm(X\tau_a(Y))-\omega_\pm(X)\,\omega_\pm(\tau_a(Y))
       \longrightarrow0\qquad(|a|\to\infty);
      \]
\item[\rm(B2)] $\omega_+\circ\theta=\omega_-$ and $\omega_-\circ\theta=\omega_+$;
\item[\rm(B3)] there are a bounded local observable $B=B^*\in\A(O)$ and $m\ne0$ such that
      $\omega_\pm(B)=\pm m$;
\item[\rm(B4)] every $\theta$-invariant ground state of $(\A,\alpha)$ equals
      $\tfrac12(\omega_++\omega_-)$.
\end{enumerate}
\end{hypothesis}

\begin{theorem}[conditional exclusion of a two-phase structure]\label{thm:dichotomy}
% src: m6/main.tex cor:nobreak (v2), statement transcribed.
Assume that $\mathscr P$ is even and $0\le\lambda/m_0^2<\varepsilon_{\mathscr P}$. Then
the assertions {\rm(B1)}--{\rm(B4)} of Hypothesis~\ref{hyp:broken} cannot hold
simultaneously.
\end{theorem}

\begin{proof}
% src: m6/main.tex cor:nobreak proof (v2), transcribed; refs adapted.
Assume that $\mathscr P$ is even and (B1)--(B4) hold. Choose $\widetilde g_n\in\G$ with
$\widetilde g_n\equiv1$ on $(-n,n)$
and write $\omega_n:=\omega_{\widetilde g_n}$. Fix $a\in\R$ such that
\[
 R_a:=\dist(O,O+a)>0,\qquad m_1R_a\ge1 ,
\]
where $O$ is the localization interval of the order parameter $B$ from (B3).

\emph{Step 1: uniform cutoff clustering for the bounded order parameter.} For the cutoff
Hamiltonian $H(\widetilde g_n)$ apply Lemma~\ref{lem:cluster} with $\gamma=m_1$
(Theorem~\ref{thm:G}), unbounded-argument slot $X=B$, bounded-argument slot
$A=\tau_a(B)\in\A(O+a)$, and localization set $\overline O$. The bounded operator $B$ has
domain $\HH$; because $B\in\A(O)\subseteq\A(U)$ for every bounded open
$U\supseteq\overline O$, it is affiliated with every such $\A(U)$, and
$\dist(\overline O,O+a)=R_a$. Thus all hypotheses of that lemma hold, and since
$\|B\Om_{\widetilde g_n}\|\le\|B\|$ and $\|\tau_a(B)\|=\|B\|$, it gives, uniformly in $n$,
\begin{equation}
 \bigl|\omega_n\bigl(B\,\tau_a(B)\bigr)-\omega_n(B)\,\omega_n(\tau_a(B))\bigr|
 \le2\|B\|^2e^{-m_1R_a/2} .
 \label{eq:cutoffcluster}
\end{equation}

\emph{Step 2: passage to the limit.} For fixed $a$, all three observables in
\eqref{eq:cutoffcluster} belong to the local algebra of one bounded interval containing
$O\cup(O+a)$. The local-norm convergence of Theorem~\ref{thm:MAIN}(1) therefore permits
passage to $n\to\infty$ in \eqref{eq:cutoffcluster}; translation invariance of
$\omega_\infty$ (Theorem~\ref{thm:MAIN}(4)) then yields
\begin{equation}
 \bigl|\omega_\infty\bigl(B\,\tau_a(B)\bigr)-\omega_\infty(B)^2\bigr|
 \le2\|B\|^2e^{-m_1R_a/2}\xrightarrow[\ |a|\to\infty\ ]{}0 .
 \label{eq:limitcluster}
\end{equation}

\emph{Step 3: the midpoint identification.} By Theorem~\ref{thm:MAIN}(3)--(4),
$\omega_\infty$ is a $\theta$-invariant ground state of $(\A,\alpha)$. Assumption (B4)
consequently gives
\[
 \omega_\infty=\tfrac12(\omega_++\omega_-),\qquad
 \omega_\infty(B)=\tfrac12(m-m)=0 ,
\]
using (B3).

\emph{Step 4: long-range order of the midpoint.} By (B1), translation invariance and
clustering of each $\omega_\pm$ imply
\[
 \omega_\pm\bigl(B\,\tau_a(B)\bigr)
 -\omega_\pm(B)\,\omega_\pm(\tau_a(B))\longrightarrow0,\qquad
 \omega_\pm(\tau_a(B))=\omega_\pm(B)=\pm m .
\]
Hence
\[
 \omega_\infty\bigl(B\,\tau_a(B)\bigr)
 =\tfrac12\!\left(\omega_+\bigl(B\,\tau_a(B)\bigr)+\omega_-\bigl(B\,\tau_a(B)\bigr)\right)
 \longrightarrow\tfrac12(m^2+m^2)=m^2\ne0 ,
\]
contradicting \eqref{eq:limitcluster}, whose left side tends to zero because
$\omega_\infty(B)=0$.
\end{proof}

\begin{remark}[scope of the phase-transition literature]\label{rem:gjs7576}
% src: lamport_gap_to_main.tex \S6; audit \S4-\S5.
The GJS phase-transition and mean-field papers \cite{GJS75,GJS76} prove the existence of
phase transitions and of two pure phases in specified strong-coupling regimes. As cited
here, they are context: they do not, by themselves, supply the full Hamiltonian
phase-space classification (B1)--(B4) --- in particular not the midpoint identification
(B4) --- for the state space of $(\A,\alpha)$. This manuscript therefore makes no
unconditional assertion that its state space carries the structure of
Hypothesis~\ref{hyp:broken}; Theorem~\ref{thm:dichotomy} says only that, at weak coupling,
that structure cannot coexist with the cutoff-uniform gap. The two rigorous options
recorded by the audit --- state (B1)--(B4) explicitly, or delete the corollary --- are
resolved here by the first.
\end{remark}

% ==================================================================
\section{Discussion}\label{sec:discussion}

% src: m6/main.tex \S8 (v2), transcribed; Part II material compressed to one paragraph (plan D8).
\emph{Relation to GRS Problems 1--2.} Problem 2 of \cite[p.~126]{GRS75} asks for local
convergence of Euclidean Gibbs states; Theorem~\ref{thm:MAIN}(1) is a Hamiltonian
counterpart, in norm on local von Neumann algebras, for the vacuum family. Problem 1 asks
for uniform local $L^p$ bounds on the Euclidean vacuum density. The note added on
\cite[p.~128]{GRS75} records Newman's solution of the small-coupling $p=1$ versions, so
the printed problems must not be read as an unqualified current open-problem claim. The
present proof neither establishes nor needs Problem 1: the identification step uses the
exact spectral identity of Proposition~\ref{prop:spectral}, and the convergence step uses
the Cauchy estimate.

\emph{Sharpness.} In the exactly solvable quadratic model the true rate is
$\mathrm{poly}(d)\,e^{-2m_1d}$ (companion note \texttt{m2/}: a round trip through the
region where the two Hamiltonians differ, at the local Agmon mass). $\Psi_{m_1}$ is
super-polynomial but sub-exponential. The present filter argument does not obtain the
sharp exponential rate; a proof of such a rate would require an additional resolvent or
Combes--Thomas input not supplied here.

\emph{Uniqueness.} $\omega_\infty$ is canonical for the cutoff family
(Corollary~\ref{cor:canonical}), and translation invariant without averaging. The theorem
does not identify every locally normal translation-invariant ground state of
$(\A,\alpha)$ with $\omega_\infty$. In particular, the Euclidean results \cite{AHZ,BDW}
do not by themselves provide the reverse reconstruction and state-space identification
needed for that Hamiltonian conclusion; companion note \texttt{m3/} records that missing
implication as its item (R3).

\emph{Beyond weak coupling.} Theorem~\ref{thm:dichotomy} is conditional on
Hypothesis~\ref{hyp:broken}. Under exactly those hypotheses, a proof based on a
cutoff-uniform gap cannot extend to that two-phase regime. The phase-transition results
\cite{GJS75,GJS76}, as cited here, are not used as a substitute for (B4)
(Remark~\ref{rem:gjs7576}), and this manuscript makes no unconditional assertion that its
Hamiltonian state space has the classification (B1)--(B4).

\emph{Part II.} Uniformity in a lattice/wavelet resolution --- the operator-algebraic
renormalization programme over the maps of \cite{MMST} --- is a separate problem with no
mathematical dependence in either direction, and its record lives in
\texttt{part\_ii/} and \texttt{strategy.md}. Its present state, in one paragraph: for the
free nearest-neighbour model, with lattice vacua $\omega^{(0)}_N$ and continuum vacuum
$\omega^{(0)}$, there are local Weyl unitaries with
$\liminf_{N\to\infty}\|\omega^{(0)}_N-\omega^{(0)}\circ\beta_N\|_{\mathfrak W_N(O)}
\ge e^{-\sqrt{2+\sqrt2}/4}-e^{-3\pi/16}>0$ for the Daubechies observable embedding
$\beta_N$, so no diagonal (finest-scale) local-state-norm error
$C\varepsilon_N^{\vartheta}$, $\vartheta>0$, is possible, and the meaningful comparison
is at fixed coarse scale; there, the free full sequence converges on each fixed coarse
Weyl generator with an explicit power rate, an interacting fixed-torus projective cluster
state exists unconditionally along a cofinal subnet, and the full-sequence identified
interacting theorem is reduced to a single scalar characteristic-functional limit, which
remains open.

\subsection*{Standing scope}

Everything in this manuscript is restricted to the weak-coupling region
$0\le\lambda/m_0^2<\varepsilon_{\mathscr P}$ of \S\ref{subsec:weak}; no claim is made for
the full noncritical region. The external residue --- a second-reader verification of the
GJS page images --- is recorded in Appendix~\ref{app:provenance}.

% ==================================================================
\appendix

\section{Constants}\label{app:constants}

All values below are fixed by the printed norms \eqref{eq:printednorm} (gate R2, page-image
reading) and the displays of the body sections; each row was re-verified against its
defining display at gates G3a/G3b (see \texttt{part\_i/README.md}).

\begin{center}
\small
\begin{tabular}{lll}
\toprule
constant & value / definition & defined in\\
\midrule
$K_1,K_2,C_7,C_8,m_1$ & GJS constants; normalization $K_2\ge1$ & Assumption~\ref{A:gjs}\\
$\varepsilon_{\mathscr P}$ & $\lambda_*/M_*^2$ & \S\ref{subsec:weak}\\
$N_I$ & $\le\lfloor L\rfloor+2\le L+2$ ($|I|=L$) & Lemma~\ref{lem:cells}\\
$\mathfrak K(p')$ & $K_1K_2^{8p'}(4p')!$ & Thm.~\ref{thm:M2}
% src: m4/r2-check.tex Prop. 2.1.
\\
$M^2$ & $9\,C_8\,\mathfrak K(p')\,\lambda^2\bigl(\sum_{n=0}^{2p'}|a_n|\bigr)^2$ & Thm.~\ref{thm:M2}\\
$M^2$ ($\varphi^4$) & $9\cdot8!\,C_8K_1K_2^{16}\lambda^2=362\,880\,C_8K_1K_2^{16}\lambda^2$ & Cor.~\ref{cor:M2phi4}\\
$K'$ & $K_2^{2}N_I^{1/2}\|f\|_2$ (exponential-moment base) & Lemma~\ref{lem:expmom}\\
$\widehat C_Q$ & $\max\{C_7e^{2m_1}\|Q\|_{\rm GJS}^2,\ e^{3m_1}\|\psi_Q\|^2\}$ & Prop.~\ref{prop:semidecay}\\
$\mathfrak m$ & total-variation bound, separation lemma & Lemma~\ref{lem:separation}\\
$d_0$ & $4+2/\gamma$ & Prop.~\ref{prop:onepath}\\
shell count & $\le8$ indices $j$ with $r_j\in[n,n+1)$ & Prop.~\ref{prop:onepath}\\
$\Psi_\gamma(d)$ & $32M\sum_{n\ge\lfloor d\rfloor}\bigl(\|W_\gamma\|_1e^{-\gamma n/4}+T_\gamma(n/2)\bigr)$; $\max\{2,\cdot\}$ if $d<d_0$ & Lemma~\ref{lem:rate}\\
$\Psi_\gamma$ tail & $\le C_{N,\gamma,M}(1+d)^{-N}$ for every $N$ & Lemma~\ref{lem:rate}\\
\bottomrule
\end{tabular}
\end{center}

\section{Provenance, corrections history, and dossier}\label{app:provenance}

\subsection*{Frozen sources of this manuscript}

\begin{center}
\small
\begin{tabular}{p{0.42\textwidth} p{0.47\textwidth}}
\toprule
file & role\\
\midrule
\texttt{m6/main.tex} (v2, audited repair) & predecessor manuscript; statements and compressed proofs\\
\texttt{audit\_2026/}\allowbreak\texttt{critical\_proof\_audit.tex} & independent audit (frozen against the v1 SHA-256)\\
\texttt{audit\_2026/}\allowbreak\texttt{lamport\_euclidean\_to\_gap.tex} & fully explicit proof, GJS $\Rightarrow$ (M$_2$)+(G)\\
\texttt{audit\_2026/}\allowbreak\texttt{lamport\_gap\_to\_main.tex} & fully explicit proof, (M$_2$)+(G) $\Rightarrow$ main theorem\\
\texttt{audit\_2026/}\allowbreak\texttt{citation\_screenshot\_dossier.tex} + \texttt{evidence/} & page images for every import\\
\texttt{m4/r1-check.tex} & gate R1: $h$-uniformity of the Euclidean inputs\\
\texttt{m4/r2-check.tex} & gate R2: printed norms and exact constants (page image, p.~594)\\
\texttt{m0/m0.tex}, \texttt{m2/m2.tex}, \texttt{m3/m3.tex} & spectral passage; quadratic-model rate; dichotomy\\
\texttt{m4/m2-audit.tex}, \texttt{m4/g-side.tex} & superseded derivations (with update notices)\\
\texttt{strategy.md} & programme record, incl.\ the unification plan and gates\\
\bottomrule
\end{tabular}
\end{center}

The two structured proof reconstructions follow the hierarchical proof format of
\cite{Lamport}.

\subsection*{Corrections history (chronological)}

\begin{itemize}
\item \emph{v0 $\to$ v1 merge} (9 Aug 2026): filter-transform sign corrected to $+1/E$;
spurious factor $\tfrac12$ removed from the one-path bound; an unnecessary $g''$ detour
eliminated; the averaging identity of the canonical-vacuum corollary repaired.
\item \emph{Gate R2} (9 Aug 2026, page-image inspection of \cite[p.~594]{GJS}): prefactor
$K_1$ to the first power; local weight $K_2^{2N(\Delta)}$; the earlier ``correction''
$K_2^{4p}\to K_2^{2p}$ of \texttt{m4/m2-audit.tex}, Rem.~5.3(i), was an OCR artifact and is
withdrawn; a half-degree slip in the v0 assembly ($(2p')!$, $K_2^{4p'}$ in place of % negative-scope mention (whitelisted): half-degree history
$(4p')!$, $K_2^{8p'}$) found and fixed. The cell-pair count for (M$_2$) is $9$, not $32$.
\item \emph{Independent audit} (9 Aug 2026): severity register
P1 (two-phase corollary missing premise (B4) --- repaired by conditioning on
Hypothesis~\ref{hyp:broken}); S1--S3 (citation pinpoints: FKN, hypercontractivity, domain
smoothing --- repaired in Assumptions \ref{A:fkn}, \ref{A:hyper}, \ref{A:domain});
L1--L6 (cell count $\lfloor L\rfloor+2$; negative-time norm ratio $c_g^{-1}$; nonnegative
FKN inputs $F_1=F_2=|V(h)|$; signed/clipped truncations; exponential-moment prefactor
$C_8K_1$; spectral-parameter range $0\le\beta<m_1$); C1 (dimensionless scaling reduction
made explicit); C2 (all imports enumerated --- Assumptions \ref{A:cutoff}--\ref{A:gjs});
H1--H2 (literature claims narrowed; GRS $p=1$ note recorded). No P0 defect was found.
\item \emph{Withdrawn routes} (kept in the record, not used here):
\texttt{m4/h6-attack.tex} Thm.~2.3 (withdrawn; UV-divergent double commutator);
gap-based (H6) proofs; wavelet-MRA-as-Polchinski decomposition (see Part II documents).
\end{itemize}

\subsection*{External residue}

Audit closure item C5 --- verification by a second human reader of the GJS page
images (1.1.7)--(1.1.10), the sharp-time localization-square convention, the scope of
$U_{\rm GJS}$, and (added 5 September 2026) the two passages of
Assumption~\ref{A:gjs}(vii$'$): the affine-path application of Theorem 1.1.7 on p.~595
and the derivation of Theorem 1.1.7 from Theorem 1.1.8 on pp.~596--597 --- remains open
at the level of this research record.  The itemized reading checklist is
\texttt{part\_i/C5\_reading\_record.md}.  The dossier and gates R1/R2 are evidence that
the imported quantifier is the source's own; they are not labelled a proof of it.

\section{$h$-uniformity: source audit and analytic prerequisites}\label{app:huniform}

% src: m4/r1-check.tex \S\S3-5, transcribed; V(h) -> \mathcal U_h, class C -> C_E, compact set K -> Lambda.
This appendix has two sharply separated purposes.  First, it records why the wording on
GJS p.~629 supports the slab-uniform interface $U_{\rm GJS}$ of
Assumption~\ref{A:gjs}(vi); the scope of that interface is fixed by the printed passages
of Assumption~\ref{A:gjs}(vii$'$) (pp.~594--597 and 629).  Second, it proves two
$h$-sensitive analytic prerequisites over the full class $\CE$: normalization of $dq_h$
and uniform $L^p$ stability under the free measure.  It does \emph{not} reproduce the
polymer expansion, its convergence summations, or the proof of all cluster-expansion
constants.  It therefore does not re-derive $U_{\rm GJS}$; that statement remains an
import of the Main Theorem, at the source's own scope.  Theorem 1.1.7 needs no separate
scope discussion because its own statement makes its constants independent of $h$ (its
printed proof consumes Theorem 1.1.8 at the same $h$).

\subsection*{Mechanism audit}

The expansion machinery of \cite[\S\S2--4]{GJS} manipulates integrals
$\int R(\Phi)\,B\,e^{-\mathcal U_h}\,d\mu_{0,C}$, where $C$ is an interpolated
covariance and $d\mu_{0,C}$ denotes the Gaussian measure with covariance $C$, through
five devices; the following lists every point of contact with $h$.
% src: m4/r1-check.tex \S3, mechanisms 1-5, transcribed.

\paragraph{M1 (the norms are kernel norms; no $h$).}
The norm (2.4.3) of an expansion term $m_\alpha$ is
$\int du\,ds\,\exp[\cdots]\,\|w(\alpha,u,s,y,\cdot)\|_{L^2}$ times combinatorial weights
(2.4.4)--(2.4.5) built from $K_1,K_2,N(\Delta)$, times $e^{\tau t}$: a functional of the
\emph{kernel} $w$ alone. No Gaussian expectation, no $e^{-\mathcal U}$, no $h$ appears in
the norm. The norm-decrease Theorems 2.4.1, 3.1, 4.1 are inequalities between such kernel
functionals, and their proofs (Lemmas 2.4.2--2.4.4, (3.2)--(3.6), \S4) estimate kernels:
contraction kernels $\delta C$ against vertex kernels.

\paragraph{M2 (covariances are $h$-free).}
The interpolation $C(s)=sC'+(1-s)C''$ (2.0.1) is between free and Dirichlet-on-contours
covariances; the kernels (2.2.2)--(2.2.3) are built from $(-d^2/dx^2+m_0^2)$ and contour
geometry. Verified against the legible scan (\texttt{m4/r1-check.tex}): no $h$ enters
any covariance, any contraction kernel, or any of the Propositions 2.1.$x$.

\paragraph{M3 ($h$ in vertices, majorised at first contact).}
Vertices arise from functional derivatives of the spectator:
\[
  \frac{\delta}{\delta\Phi(x)}e^{-\mathcal U_h}
  =-\lambda\,h(x)\,{:}\mathscr P'(\Phi(x)){:}\;e^{-\mathcal U_h},
\]
and iterates with $\mathscr P'',\dots$ A vertex therefore contributes to the \emph{kernel}
a factor $\lambda h(x)p_j$ supported in a lattice square. The very first estimate applied
to it is the printed one (p.~626): ``$v(x)$ is $O(\lambda)$ times the characteristic
function of a lattice square'' --- that is, $|\lambda h(x)p_j|\le O(\lambda)\chi_\Delta(x)$,
using $0\le h\le1$ and nothing else. Downstream of this line the kernels are $h$-free.
This is the single point where $h$ touches a quantitative estimate.

\paragraph{M4 (the identities are exact and $h$-blind).}
The devices (2.0.4)--(2.0.5) (covariance interpolation), (2.2.1), (2.2.5), (2.2.7)
(advancing annihilators), and (4.1)--(4.2), (4.5) (integration by parts) are identities
applied to integrands $F=R(\Phi)B_Te^{-\mathcal U_h}$; $e^{-\mathcal U_h}$ rides along
inside $F$. Their validity requires only that the integrals exist --- integrability, not
bounds. For the free-measure endpoint this is supplied, uniformly over $\CE$, by
Theorem~\ref{thm:stability}; existence under the interpolated measures $d\mu_{0,C}$ is
internal to the printed proof and is consumed here only as such (by M5, no constant of
the source depends on it).

\paragraph{M5 (termination cancels every $e^{-\mathcal U}$; no stability constant).}
The full expansion (4.3) reads
\[
  \int AB_T\,e^{-\mathcal U}d\mu_{0,C_0}
  =\sum_i\int\Bigl[\int dy\,w_i(y)\,{:}\Phi(y_1)\cdots\Phi(y_i){:}\Bigr]B_T\,
   e^{-\mathcal U}d\mu_{0,C_0}
  \;+\;\sum_r w_r\int e^{-\mathcal U}d\mu_{0,C_0},
\]
with $w_i,w_r$ \emph{numerical kernels} controlled by the norm inequalities; the
remainder terms are literally multiples of $\int e^{-\mathcal U}d\mu_0$ (their Table~1:
final-$r$ terms have ``no vertices, no annihilators''). The proof of Theorems
1.1.8/1.1.11 then \emph{divides (4.3) by $\int e^{-\mathcal U}d\mu_{0,C_0}$} (p.~628):
the first sum reassembles into $dq_h$-integrals of the comparison polynomials, and the
remainder becomes $\sum_rw_r$ exactly. Every factor $e^{-\mathcal U_h}$ cancels against
the normalisation. Consequently no upper or lower bound on
$\int e^{-\mathcal U}d\mu_0$, and no stability estimate, enters any constant of the
proof.

\begin{proposition}[consequence of the audited source mechanisms]\label{prop:struct}
% src: m4/r1-check.tex prop:struct, transcribed.
Within mechanisms M1--M5 as recorded above, every displayed quantitative bound is
unchanged when $h$ ranges over $0\le h\le1$: M3 replaces its only vertex occurrence by
$|h|\le1$, while M4--M5 use identities or exact cancellation.  Consequently the
constants visible in those mechanisms depend on
$(m_0,K_2,K_1,\lambda,n,\varepsilon,\mathscr P)$ and not on $h$.
\end{proposition}

\begin{proof}
By direct substitution in M1--M5: the quantitative chain displayed there is entirely at the level of kernel norms (M1), whose
ingredients are $h$-free covariance/contraction kernels (M2) and vertex kernels majorised
by $O(\lambda)\chi_\Delta$ (M3); the measure-level steps are exact identities (M4) or
exact cancellations (M5).  This proposition is deliberately limited to the audited
mechanisms; it makes no completeness claim about portions of the source expansion not
reproduced here.
\end{proof}

\subsection*{The two $h$-sensitive analytic facts, proved}

Notation for this subsection: $G:=(-\Delta_{\R^2}+m_0^2)^{-1}$ is the Euclidean
covariance, $G(\zeta)=\frac1{2\pi}K_0(m_0|\zeta|)$; $h\in\CE$ with
$\supp h\subseteq\Lambda$, $\Lambda$ compact;
$\mathcal U_h=\lambda\int h\,{:}\mathscr P(\Phi){:}\,dz$ and
$Z(h)=\int e^{-\mathcal U_h}d\mu_0$ as in Assumption~\ref{A:stability}. Choose once and
for all
\[
 \delta\in C_c^\infty(\R^2),\qquad
 \delta\ge0,\qquad \delta(z)=\delta(-z),\qquad
 \supp\delta\subseteq\{z:|z|\le1\},\qquad
 \int_{\R^2}\delta(z)\,dz=1,
\]
and, for $\kappa\ge1$, define the two-dimensional rescaling
$\delta_\kappa(z):=\kappa^2\delta(\kappa z)$.  Thus convolution by
$\delta_\kappa$ is a positive average over a ball of radius $\kappa^{-1}$; a double
convolution averages over displacements of radius at most $2\kappa^{-1}$.  Set
$\Phi_\kappa=\delta_\kappa*\Phi$,
$c_\kappa(z)=\E[\Phi_\kappa(z)^2]\le\frac1{2\pi}\log\kappa+C_*$ (translation invariance),
and $\mathcal U_{h,\kappa}=\lambda\int h\,{:}\mathscr P(\Phi_\kappa){:}_{c_\kappa}dz$,
Wick ordered with respect to the variance of the cutoff field, so that
$\mathcal U_{h,\kappa}\to\mathcal U_h$ in every $L^p(d\mu_0)$.

\begin{lemma}[normalisation]\label{lem:Z}
% src: m4/r1-check.tex lem:Z, transcribed.
For every $h\in\CE$ with $\supp h\subseteq\Lambda$:
$Z(h)\ge e^{-\lambda|a_0|\,|\Lambda|}>0$, and $Z(h)<\infty$ by
Theorem~\ref{thm:stability}. Hence $dq_h$ is a well-defined probability measure on the
whole class $\CE$.
\end{lemma}

\begin{proof}
$\E[\mathcal U_h]=\lambda a_0\int h\in[-\lambda|a_0||\Lambda|,\lambda|a_0||\Lambda|]$,
since Wick monomials of degree $\ge1$ have mean zero and
$0\le h\le\mathbf 1_\Lambda$. Jensen's inequality for the probability measure $d\mu_0$
gives $Z(h)\ge e^{-\E[\mathcal U_h]}\ge e^{-\lambda|a_0||\Lambda|}$.
\end{proof}

\begin{lemma}[Wick lower bound; the $h\ge0$ entry]\label{lem:wick}
% src: m4/r1-check.tex lem:wick, transcribed.
There is $b=b(\mathscr P)$ such that for all $c\ge0$ and $x\in\R$,
${:}\mathscr P{:}_c(x):=\sum_na_nc^{n/2}\operatorname{He}_n(x/\sqrt c)\ge-b\,(1+c)^{p'}$.
Consequently, for every $h\in\CE$ with $\supp h\subseteq\Lambda$ and every
$\kappa\ge2$,
\[
  \mathcal U_{h,\kappa}\ \ge\
  -\lambda\,b\,|\Lambda|\,\bigl(1+\tfrac1{2\pi}\log\kappa+C_*\bigr)^{p'}
  \ \ge\ -c_1(\lambda,\mathscr P,\Lambda)\,(\log\kappa)^{p'} .
\]
\end{lemma}

\begin{proof}
Expanding the Hermite polynomials,
${:}\mathscr P{:}_c(x)=\sum_{m=0}^{2p'}\beta_mx^m$ with $\beta_{2p'}=a_{2p'}$ and
$|\beta_m|\le C(\mathscr P)(1+c)^{(2p'-m)/2}$ (each power $c^j$ comes with $x^{n-2j}$).
For a polynomial $Ax^{2p'}+\sum_{m<2p'}\beta_mx^m$ with $A>0$, weighted AM--GM gives, for
each $m<2p'$,
\[
  |\beta_m||x|^m\ \le\ \frac{A}{4p'}\,x^{2p'}
  +C_{p'}\Bigl(\frac{|\beta_m|^{2p'}}{A^{m}}\Bigr)^{\!1/(2p'-m)},
\]
and summing,
$\min_x{:}\mathscr P{:}_c\ge-C_{p'}'\sum_m(|\beta_m|^{2p'}/A^m)^{1/(2p'-m)}$. With
$|\beta_m|\le C(1+c)^{(2p'-m)/2}$ the exponent of $(1+c)$ is
$\frac{2p'}{2p'-m}\cdot\frac{2p'-m}{2}=p'$, uniformly in $m$. The second claim follows
from $h\ge0$, $\int h\le|\Lambda|$, and the pointwise bound applied under the integral
--- this is the single use of the \emph{sign} of $h$.
\end{proof}

\begin{lemma}[ultraviolet difference; the $h\le1$ entry]\label{lem:uv}
% src: m4/r1-check.tex lem:uv, transcribed.
There is $c_2=c_2(\lambda,\mathscr P,\Lambda)$ such that for all $h\in\CE$ with
$\supp h\subseteq\Lambda$ and all $\kappa\ge2$,
\[
  \bigl\|\mathcal U_h-\mathcal U_{h,\kappa}\bigr\|_{L^2(d\mu_0)}^2
  \ \le\ c_2\,\kappa^{-1}(\log\kappa)^{2p'} .
\]
\end{lemma}

\begin{proof}
By the Wick isometry, with $G^{(1)}_\kappa:=(\delta_\kappa\otimes1)*G$ and
$G_\kappa:=(\delta_\kappa\otimes\delta_\kappa)*G$,
\[
  \|\mathcal U_h-\mathcal U_{h,\kappa}\|_{L^2}^2
  =\lambda^2\sum_{n=1}^{2p'}a_n^2\,n!\iint h(z)h(z')\,
  \bigl[G^n-2(G^{(1)}_\kappa)^n+G_\kappa^n\bigr](z-z')\,dz\,dz'
\]
(the $c_\kappa$-Wick ordering makes the diagonal subtractions match; degree $0$
cancels). Using $|h|\le\mathbf 1_\Lambda$ --- the single use of the \emph{upper} bound on
$h$ --- it suffices to bound $\iint_{\Lambda\times\Lambda}|G^n-J^n|$ for
$J\in\{G^{(1)}_\kappa,G_\kappa\}$. The classical Bessel facts
$0<G(\zeta)\le C(1+|\log|\zeta||)e^{-m_0|\zeta|}$ and
$|\nabla G(\zeta)|\le C(|\zeta|^{-1}+1)e^{-m_0|\zeta|}$ give, for the average $J$ of $G$
over balls of radius $\le2/\kappa$: $J(\zeta)\le C(1+\log\kappa)$ for all $\zeta$, and
$|J(\zeta)-G(\zeta)|\le C\kappa^{-1}(|\zeta|^{-1}+1)$ for $|\zeta|\ge4/\kappa$. Split the
integral over $\zeta=z-z'$:
\[
  \int_{|\zeta|<4/\kappa}\bigl[G^n+J^n\bigr]
  \le C_n\Bigl[\int_{|\zeta|<4/\kappa}(1+|\log|\zeta||)^nd\zeta
     +\kappa^{-2}(1+\log\kappa)^n\Bigr]
  \le C_n'\kappa^{-2}(\log\kappa)^n,
\]
\[
  \int_{4/\kappa\le|\zeta|\le R}n\max(G,J)^{n-1}|G-J|
  \le C_n\kappa^{-1}\!\int_0^R(1+|\log r|)^{n-1}(1+r)\,dr
  \le C_n(R)\,\kappa^{-1},
\]
with $R=\operatorname{diam}\Lambda$; multiply by $|\Lambda|$ for the outer variable and
sum over $n\le2p'$.
\end{proof}

\begin{theorem}[$h$-uniform stability]\label{thm:stability}
% src: m4/r1-check.tex thm:stability, transcribed.
For every $p<\infty$ and compact $\Lambda$,
\[
  \sup\bigl\{\|e^{-\mathcal U_h}\|_{L^p(d\mu_0)}\ :\ h\in\CE,\
  \supp h\subseteq\Lambda\bigr\}\ <\ \infty .
\]
\end{theorem}

\begin{proof}
(For $\lambda=0$, $\mathcal U_h=0$ and the statement is trivial; assume $\lambda>0$.)
$Y:=\mathcal U_h-\mathcal U_{h,\kappa}$ is a sum $Y=\sum_{r=0}^{2p'}Y_r$ of Wick chaoses,
so Nelson's hypercontractive moment comparison applies in the form
\[
 \|Y\|_{L^q}\le(q-1)^{p'}\|Y\|_{L^2}\qquad(q\ge2) .
\]
(Proof from the abstract form of Assumption~\ref{A:hyper} \cite[Thm.~I.17]{SimonBook}:
the second-quantized contraction $\Gamma(e^{-t})$ on $L^2(d\mu_0)$ satisfies
$\|\Gamma(e^{-t})X\|_q\le\|X\|_2$ for $q=1+e^{2t}$, and
$\Gamma(e^{-t})Y_r=e^{-rt}Y_r$. Write $Y=\Gamma(e^{-t})Z$ with $Z=\sum_re^{rt}Y_r$; by
orthogonality of the chaoses, $\|Z\|_2\le e^{2p't}\|Y\|_2$, whence
$\|Y\|_q\le\|Z\|_2\le e^{2p't}\|Y\|_2=(q-1)^{p'}\|Y\|_2$.)

For $b\ge b_0(\lambda,\mathscr P,\Lambda)$ choose $\kappa=\kappa(b)$ by
$c_1(\log\kappa)^{p'}=b/2$, i.e.\ $L:=\log\kappa=(b/2c_1)^{1/p'}$. By
Lemma~\ref{lem:wick},
$\{\mathcal U_h\le-b\}\subseteq\{\mathcal U_h-\mathcal U_{h,\kappa}\le-b/2\}$, so by
Chebyshev in $L^{2q}$ and Lemma~\ref{lem:uv},
\[
  \Pr\bigl[\mathcal U_h\le-b\bigr]
  \ \le\ \Bigl(\frac2b\Bigr)^{2q}(2q)^{2p'q}\,
  \bigl(c_2\,e^{-L}L^{2p'}\bigr)^{q}
  \ =\ \Bigl[\frac{2\sqrt{c_2}}{b}\,(2q)^{p'}\,L^{p'}e^{-L/2}\Bigr]^{2q}.
\]
Choose $2q=\exp\bigl(\frac{L}{4p'}\bigr)$, so that $(2q)^{p'}=e^{L/4}$ and the bracket is
$\le\frac{2\sqrt{c_2}}{b}L^{p'}e^{-L/4}\le e^{-1}$ for
$b\ge b_1(\lambda,\mathscr P,\Lambda)$, whence
\[
  \Pr\bigl[\mathcal U_h\le-b\bigr]\ \le\
  \exp\Bigl(-\exp\Bigl(\tfrac1{4p'}(b/2c_1)^{1/p'}\Bigr)\Bigr),
\]
doubly exponentially small, uniformly over the stated class; and
$\exp\bigl(\tfrac1{4p'}(b/2c_1)^{1/p'}\bigr)\ge2pb$ for $b$ large, since a power of $b$
inside an exponential beats $\log b$. Then
$\|e^{-\mathcal U_h}\|_{L^p}^p
=\int_0^\infty pe^{pb}\Pr[\mathcal U_h\le-b]\,db+O(e^{pb_1})<\infty$, with a bound
depending only on $(p,\lambda,\mathscr P,\Lambda)$. Every constant above came from
Lemmas~\ref{lem:wick}--\ref{lem:uv}, which used exactly $h\ge0$ and $h\le1$.
\end{proof}

\begin{remark}[why the class is what it is]\label{rem:whyclass}
% src: m4/r1-check.tex, closing remark of \S4, transcribed.
The two halves of the constraint $0\le h\le1$ are used in complementary places:
positivity in the pointwise lower bound (Lemma~\ref{lem:wick} would fail for $h$ of
variable sign --- the interaction density is unbounded above \emph{and} below), and
boundedness in the $L^2$ kernel estimates (Lemma~\ref{lem:uv} and M3). Nothing else about
$h$ --- smoothness, shape, time extent, monotonicity along a sequence --- is touched
anywhere. This is the structural content of Source Fact~\ref{sf:printed}.
\end{remark}

\begin{proposition}[coverage of every CE application]\label{prop:discharge}
% src: m4/r1-check.tex thm:verdict, transcribed.
Assume $U_{\rm GJS}$.  Then every use of CE1 in
\S\S\ref{sec:slab}--\ref{sec:G} has a cutoff in $\CE^{\rm slab}$ and is covered by
\eqref{eq:UGJS}; every use of CE2 is covered directly by Source
Fact~\ref{sf:CE2}, whose constants are independent of $h$.  Lemma~\ref{lem:Z} and
Theorem~\ref{thm:stability} supply the normalization and integrability needed to define
the same measures.  Hence no stronger cutoff-uniformity assertion is consumed by the
proof.
\end{proposition}

\begin{proof}
The CE1 occurrences are the displays in Theorem~\ref{thm:M2},
Lemmas~\ref{lem:transfer}--\ref{lem:twopoint}, and their immediate consequences; in
each case the measure is explicitly $dq_{h_{T,g}}$.  If $g$ is replaced by an affine
path cutoff $g_s$, Lemma~\ref{lem:classes} gives $g_s\in\CH$, so
$h_{T,g_s}\in\CE^{\rm slab}$.  Lemma~\ref{lem:ce2strip} is the only connected
CE2 application and Source Fact~\ref{sf:CE2} quantifies over its cutoff directly.
The remaining assertions are exactly Lemma~\ref{lem:Z} and
Theorem~\ref{thm:stability}.
\end{proof}

% ==================================================================
\begin{thebibliography}{99}

\bibitem{GJ1} J.~Glimm, A.~Jaffe,
\emph{A $\lambda\phi^4$ quantum field theory without cutoffs. I},
Phys.\ Rev.\ \textbf{176} (1968) 1945--1951.

\bibitem{GJ2} J.~Glimm, A.~Jaffe,
\emph{The $\lambda(\varphi^4)_2$ quantum field theory without cutoffs. II. The field
operators and the approximate vacuum},
Ann.\ of Math.\ \textbf{91} (1970) 362--401.

\bibitem{GJ3} J.~Glimm, A.~Jaffe,
\emph{The $\lambda(\varphi^4)_2$ quantum field theory without cutoffs. III. The physical
vacuum},
Acta Math.\ \textbf{125} (1970) 203--267.

\bibitem{GJ4} J.~Glimm, A.~Jaffe,
\emph{The $\lambda(\varphi^4)_2$ quantum field theory without cutoffs. IV. Perturbations of
the Hamiltonian},
J.\ Math.\ Phys.\ \textbf{13} (1972) 1568--1584.

\bibitem{GJ71} J.~Glimm, A.~Jaffe,
\emph{Positivity and self-adjointness of the $P(\phi)_2$ Hamiltonian},
Comm.\ Math.\ Phys.\ \textbf{22} (1971) 253--258.

\bibitem{GJExp} J.~Glimm, A.~Jaffe,
\emph{Quantum Field Theory and Statistical Mechanics: Expositions},
Birkh\"auser, Boston, 1985.

\bibitem{GJS} J.~Glimm, A.~Jaffe, T.~Spencer,
\emph{The Wightman axioms and particle structure in the $\mathscr P(\phi)_2$ quantum field
model},
Ann.\ of Math.\ \textbf{100} (1974) 585--632.

\bibitem{GJS75} J.~Glimm, A.~Jaffe, T.~Spencer,
\emph{Phase transitions for $\varphi^4_2$ quantum fields},
Comm.\ Math.\ Phys.\ \textbf{45} (1975) 203--216.

\bibitem{GJS76} J.~Glimm, A.~Jaffe, T.~Spencer,
\emph{A convergent expansion about mean field theory. I, II},
Ann.\ Physics \textbf{101} (1976) 610--630, 631--669.

\bibitem{GRS75} F.~Guerra, L.~Rosen, B.~Simon,
\emph{The $P(\phi)_2$ Euclidean quantum field theory as classical statistical mechanics},
Ann.\ of Math.\ \textbf{101} (1975) 111--189, 191--259.

\bibitem{Rosen70} L.~Rosen,
\emph{A $\lambda\phi^{2n}$ field theory without cutoffs},
Comm.\ Math.\ Phys.\ \textbf{16} (1970) 157--183.

\bibitem{SimonBook} B.~Simon,
\emph{The $P(\Phi)_2$ Euclidean (Quantum) Field Theory},
Princeton University Press, Princeton, 1974.

\bibitem{AHZ} S.~Albeverio, R.~H\o egh-Krohn, B.~Zegarlinski,
\emph{Uniqueness and global Markov property for Euclidean fields: the case of general
polynomial interactions},
Comm.\ Math.\ Phys.\ \textbf{123} (1989) 377--424.

\bibitem{BDW} R.~Bauerschmidt, B.~Dagallier, H.~Weber,
\emph{Holley--Stroock uniqueness method for the $\varphi^4_2$ dynamics},
arXiv:2504.08606 (2025).

\bibitem{RZT} M.~Rodriguez Zarate, T.~Thiemann,
\emph{Hamiltonian renormalisation VIII. $P(\Phi)_2$ quantum field theory},
arXiv:2505.13030 (2025).

\bibitem{Hastings} M.~B.~Hastings,
\emph{Lieb--Schultz--Mattis in higher dimensions},
Phys.\ Rev.\ B \textbf{69} (2004) 104431.

\bibitem{BMNS} S.~Bachmann, S.~Michalakis, B.~Nachtergaele, R.~Sims,
\emph{Automorphic equivalence within gapped phases of quantum lattice systems},
Comm.\ Math.\ Phys.\ \textbf{309} (2012) 835--871.

\bibitem{MMST} V.~Morinelli, G.~Morsella, A.~Stottmeister, Y.~Tanimoto,
\emph{Scaling limits of lattice quantum fields by wavelets},
Comm.\ Math.\ Phys.\ \textbf{387} (2021) 1299--1367.

\bibitem{Lamport} L.~Lamport,
\emph{How to write a 21st century proof},
J.\ Fixed Point Theory Appl.\ \textbf{11} (2012) 43--63.

\end{thebibliography}

\end{document}
